\documentclass[12pt,a4paper]{ctexart}

% Packages
\usepackage{amsmath,amssymb,amsthm}
\usepackage{booktabs}
\usepackage{hyperref}
\usepackage[margin=2.5cm]{geometry}
\usepackage{setspace}
\usepackage{enumitem}
\onehalfspacing

% Theorem environments
\newtheorem{theorem}{Theorem}[section]
\theoremstyle{definition}
\newtheorem{definition}{Definition}[section]
\theoremstyle{remark}
\newtheorem{remark}{Remark}[section]

\title{The Yajie Protocol: Technology Lock-in, Audit Sovereignty,\\and the Non-Proliferation Logic of Data Quality Assessment}
\author{SCX Research Group}
\date{June 2026}

\begin{document}

\maketitle

\begin{abstract}
As artificial intelligence systems become pervasive in high-stakes decision-making, the quality of training and inference data has emerged as a critical—yet systematically under-governed—input. This paper introduces the Yajie Protocol, a multi-expert consistency-based data auditing algorithm, and analyzes its structural implications for the emerging market in data quality assessment. We demonstrate that Yajie generates a distinctive form of technology lock-in driven not by network effects in the conventional sense, nor by proprietary control of intellectual property, but by a time-accumulated advantage mechanism: each audit cycle enriches a shared evidentiary corpus, widening the quality gap between the incumbent protocol and any potential entrant. Drawing on game-theoretic modeling, we formalize this dynamic as a \textit{Non-Proliferation Equilibrium} (NPE), wherein rational actors forgo the development of competing audit standards despite the absence of formal barriers to entry. We establish a structural analogy between this equilibrium and the strategic logic underpinning the Nuclear Non-Proliferation Treaty (NPT), arguing that audit standards fragmentation produces negative externalities analogous to nuclear proliferation. The paper further examines the geopolitical dimensions of audit infrastructure through the lens of the SWIFT financial messaging system and the International Accounting Standards Board (IASB), with particular attention to US-China technological decoupling. We model the lock-in trajectory through a five-stage progression and derive policy recommendations for embedding audit neutrality within multilateral governance frameworks. Our findings carry implications for the design of international AI governance institutions, the regulation of algorithmic auditing markets, and the preservation of epistemic sovereignty in an era of data-driven geopolitical competition.

\bigskip
\noindent\textbf{Keywords:} data quality assessment, algorithmic audit, technology lock-in, non-proliferation, game theory, audit sovereignty, AI governance, multi-expert consistency, SCX uncertainty principle, mutual deterrence, adoption-audit equilibrium
\end{abstract}

\section{Introduction}

The governance of artificial intelligence has entered a phase of institutional crystallization. Where the preceding decade was characterized by proliferation of ethical principles and voluntary guidelines (Floridi et al., 2018; Jobin et al., 2019), the current moment is defined by the translation of normative commitments into operational mechanisms: conformity assessments, third-party audits, and certification regimes (Raji et al., 2020; M\"{o}kander et al., 2021). Central to this operational turn is the question of \textit{data quality}—the accuracy, completeness, consistency, and fitness-for-purpose of the datasets that underpin machine learning pipelines.

The dependency of AI system performance on data quality is by now well-established in both the technical literature (Sambasivan et al., 2021; Northcutt et al., 2021) and policy discourse (European Commission, 2021; NIST, 2023). Yet the institutional infrastructure for assessing and certifying data quality remains nascent. A small number of commercial auditing services have emerged, alongside standards development efforts within ISO/IEC JTC 1/SC 42 and the IEEE P2863 working group. What has not yet materialized—and what this paper argues may \textit{never} materialize in a competitive form—is a pluralistic market of interoperable data quality assessment protocols.

The central claim of this paper is that data quality auditing exhibits structural properties that drive the market toward a natural monopoly, and that this monopolistic tendency is, counterintuitively, \textit{welfare-enhancing} under a specific set of conditions. We develop this argument through a detailed analysis of the Yajie Protocol, a data auditing algorithm that operationalizes multi-expert consistency as a quality signal, and we formalize the resulting market dynamics through the concept of a \textit{Non-Proliferation Equilibrium} (NPE).

The Yajie Protocol, which we describe in Section~3, operates by deploying an ensemble of independent audit modules against a target dataset and measuring the degree of inter-auditor agreement. Disagreement signals potential quality defects; convergence across auditors with orthogonal detection strategies provides evidence of quality. What makes the protocol structurally distinctive is its \textit{cumulative evidentiary architecture}: every audit cycle enriches a shared repository of detected anomalies, auditor calibration parameters, and dataset-specific quality profiles, such that the protocol becomes monotonically more accurate over time. New entrants cannot replicate this accumulated capital; they must either license access to the incumbent's evidentiary corpus (entrenching its position) or enter the market with a structurally inferior product.

This dynamic is not adequately captured by standard models of technology lock-in. The QWERTY literature (David, 1985; Arthur, 1989) emphasizes network effects, switching costs, and installed-base advantages. Platform economics (Rochet \& Tirole, 2003; Parker \& Van Alstyne, 2005) highlights cross-side network externalities. Neither framework fully accounts for a product whose quality is an increasing function of \textit{cumulative usage time}, independent of the number of users. The Yajie Protocol's lock-in mechanism is temporal, not networked: it accrues advantage through the sheer volume of audits conducted, each of which makes the protocol better at conducting the next audit. This \textit{quality ratchet} creates an entry barrier that grows with time, asymptotically approaching an insurmountable moat.

We deploy a game-theoretic framework to demonstrate that, under plausible parameterizations, rational competitors will choose not to invest in the development of competing audit protocols, even when the incumbent has not engaged in predatory pricing or exclusionary practices. The resulting market structure—a single dominant protocol, universally adopted—resembles not a coercive monopoly of the Standard Oil variety, but rather the equilibrium that obtains in domains where fragmentation itself is the primary source of social cost. We further demonstrate that the protocol's architecture generates a distinctive \textit{Adoption-Audit Equilibrium}: organizations are simultaneously incentivized to adopt Yajie (because it improves data quality) and deterred from falsifying data (because adoption renders audit outputs unconditionally verifiable by third parties, per the SCX Uncertainty Principle). The architecture transforms data quality from an unverifiable assertion into a revealed attribute—and honesty into the dominant strategy.

This is where the nuclear non-proliferation analogy becomes analytically productive. The Nuclear Non-Proliferation Treaty (NPT) succeeded not because it prohibited nuclear weapons development (it did not, for the nuclear-weapon states) but because it established a club whose benefits—access to peaceful nuclear technology, participation in the global non-proliferation regime—outweighed the strategic value of an independent nuclear program for most states (Sagan, 1996). Critically, the NPT's stability does not depend on enforcement alone; it depends on a rational calculation by non-nuclear states that proliferation would leave them \textit{worse off} than the status quo, either because their nascent capability would be neutralized by preemptive action, because the resultant arms race would consume resources better spent elsewhere, or because proliferation by one state would trigger proliferation by its rivals, leaving everyone less secure (Waltz, 1981; Mearsheimer, 1993).

The audit protocol market exhibits an analogous structure. A world with $N$ audit protocols, each partially trusted and partially contested, is a world in which audit results are themselves subject to meta-audit disputes, eroding the very certainty that audits are designed to provide. The development of a competing protocol by one jurisdiction would trigger reciprocal development by rivals, fragmenting the global audit infrastructure along geopolitical lines and imposing coordination costs on every entity that operates across jurisdictions. In this environment, the \textit{rational} choice for most actors is to accept the dominance of a single protocol—provided that protocol is governed under adequate neutrality safeguards—rather than to pursue audit autarky.

The paper proceeds as follows. Section~2 provides a technical overview of the Yajie Protocol and situates it within the literature on algorithmic auditing and technology lock-in. Section~3 develops the Non-Proliferation Equilibrium model, including the formal payoff structure, Nash equilibrium characterization, and NPT analogy. Section~4 examines the geopolitical dimensions of audit infrastructure through case studies of SWIFT and the IASB, with an extension to the US-China technology competition. Section~5 models the lock-in trajectory through five stages and analyzes the business model dynamics. Section~6 derives policy recommendations for embedding audit neutrality within multilateral governance. Section~7 addresses limitations and identifies directions for future research. Section~8 concludes. Section~9 synthesizes the foregoing analyses into a unified game-theoretic architecture.

\section{Theoretical Background}

\subsection{The Yajie Protocol: A Technical Overview}

The Yajie Protocol is a computational framework for assessing the quality of structured and semi-structured datasets through multi-expert consistency analysis. At its core, Yajie operationalizes a simple but powerful principle: when multiple independent audit procedures, each designed to detect distinct categories of data defect, converge on a quality assessment, the probability that the assessment is correct increases multiplicatively. Conversely, when auditors diverge, the pattern of divergence carries diagnostic information about the nature and location of quality failures.

Formally, let $\mathcal{D}$ be a target dataset consisting of $n$ records, each characterized by a feature vector $\mathbf{x}_i \in \mathbb{R}^d$ and, where applicable, a label $y_i$. Let $\mathcal{A} = \{A_1, A_2, \ldots, A_k\}$ be an ensemble of $k$ audit modules, where each $A_j: \mathcal{D} \rightarrow \mathcal{R}_j$ maps the dataset to a structured audit report $\mathcal{R}_j = (\mathbf{q}_j, \mathbf{a}_j, \mathbf{c}_j)$ comprising:
\begin{itemize}
\item A quality score vector $\mathbf{q}_j \in [0,1]^m$ across $m$ quality dimensions (e.g., completeness, accuracy, consistency, timeliness, uniqueness);
\item An anomaly register $\mathbf{a}_j = \{(i, \tau, s)\}$ recording the index $i$ of each flagged record, the anomaly type $\tau$, and a severity score $s \in [0,1]$;
\item A confidence vector $\mathbf{c}_j \in [0,1]^m$ expressing the auditor's self-assessed reliability on each dimension.
\end{itemize}

The protocol's aggregation function $\Phi: \prod_{j=1}^k \mathcal{R}_j \rightarrow \mathcal{R}^*$ synthesizes the ensemble outputs into a unified audit report $\mathcal{R}^*$ through a weighted consensus mechanism:
\begin{equation}
\Phi(\mathcal{R}_1, \ldots, \mathcal{R}_k) = \left( \sum_{j=1}^k w_j \mathbf{q}_j, \; \bigcup_{j=1}^k \mathbf{a}_j \setminus \mathcal{F}, \; \prod_{j=1}^k \text{diag}(\mathbf{c}_j) \right)
\end{equation}
where $w_j$ are dynamic weights derived from each auditor's historical calibration performance, $\mathcal{F}$ is a false-positive filter trained on previously adjudicated anomalies, and the confidence aggregation employs element-wise multiplication to reflect the multiplicative reduction in uncertainty when independent auditors concur.

The protocol's defining architectural feature is the \textit{Cumulative Evidentiary Corpus} (CEC), denoted $\mathcal{E}_t$, which aggregates all audit outputs up to time $t$:
\begin{equation}
\mathcal{E}_t = \bigcup_{\tau=0}^{t} \{\mathcal{R}^*_\tau, \mathcal{J}_\tau, \Delta\mathcal{P}_\tau\}
\end{equation}
where $\mathcal{J}_\tau$ is the set of human-adjudicated judgments from audit cycle $\tau$, and $\Delta\mathcal{P}_\tau$ represents parameter updates derived from those judgments. The CEC serves three functions: (i) it provides training data for the false-positive filter $\mathcal{F}$; (ii) it enables dynamic recalibration of auditor weights $w_j$; and (iii) it accumulates a growing library of dataset-specific quality fingerprints that accelerate future audits of similar or related datasets.

The CEC grows with each audit cycle. Let $g(t)$ denote the cumulative number of audits conducted by time $t$. Each audit contributes, on average, $\bar{\alpha}$ new anomaly instances and $\bar{\beta}$ new calibration data points to the CEC. The protocol's accuracy $\theta(t)$—measured as the F1 score of its anomaly detection relative to ground truth—improves as a function of CEC size:
\begin{equation}
\theta(t) = \theta_{\max} - (\theta_{\max} - \theta_0) e^{-\gamma \cdot |\mathcal{E}_t|}
\end{equation}
where $\theta_{\max}$ is the asymptotic accuracy ceiling, $\theta_0$ is the initial accuracy, and $\gamma$ is a learning rate parameter. This functional form captures diminishing returns to additional data, consistent with the empirical scaling literature (Kaplan et al., 2020), while ensuring that the incumbent's accuracy advantage over a zero-CEC entrant grows monotonically with cumulative audit volume.

\subsection{Technology Lock-in: Classical Mechanisms and Their Limits}

The phenomenon of technology lock-in—whereby an inferior technology persists because the costs of coordinated transition exceed the benefits of switching to a superior alternative—has been a central concern of innovation economics since David's (1985) iconic analysis of the QWERTY keyboard. Arthur (1989) formalized the mechanisms through which increasing returns to adoption generate path-dependent outcomes: large fixed costs, learning effects, coordination effects, and adaptive expectations.

Subsequent literature refined and extended these mechanisms. Katz and Shapiro (1985, 1994) developed the economics of network effects, distinguishing between direct network effects (the value of a network increases with the number of users, as in telecommunications) and indirect network effects (a larger installed base attracts complementary goods, as in software platforms). Farrell and Saloner (1985) modeled the dynamics of standardization, demonstrating that excess inertia—the failure to adopt a superior standard because of coordination failure—can arise even when all actors prefer the superior standard.

The two-sided market literature (Rochet \& Tirole, 2003; Armstrong, 2006) extended these insights to platforms that intermediate between distinct user groups, showing that cross-side network externalities can generate powerful lock-in effects even in the absence of traditional switching costs. Once a platform achieves critical mass on both sides of the market, entrants face a chicken-and-egg problem: they cannot attract users from one side without first assembling a critical mass on the other.

These classical mechanisms do not, however, capture the lock-in dynamic generated by the Yajie Protocol. Consider the disanalogies:

\begin{enumerate}[label=\arabic*.]
\item \textbf{Network effects}: Yajie's value does not depend on the number of other users of the protocol. A single organization conducting repeated audits accrues the same CEC-driven accuracy gains regardless of whether other organizations also use the protocol. The lock-in is \textit{temporal}, not \textit{networked}.

\item \textbf{Switching costs}: While Yajie does generate switching costs in the conventional sense (retraining staff, reconfiguring pipelines, renegotiating contracts), the principal entry barrier is not the cost of switching \textit{away} from Yajie but the cost of replicating the \textit{accumulated audit history} that makes Yajie accurate. A competitor cannot buy, license, or reverse-engineer the CEC, because the CEC is not merely a database but a cumulative record of audit outcomes, adjudications, and parameter updates that are causally entangled with the protocol's specific architecture.

\item \textbf{Installed base}: Unlike a platform whose value derives from its user base, Yajie's value derives from its \textit{audit base}—the corpus of audits it has conducted. This is a stock variable, not a flow variable, and it cannot be replicated by attracting users away from the incumbent.

\item \textbf{Complementary goods}: The Yajie ecosystem does generate complementary assets—training programs, integration tooling, certification bodies—but these are derivative of the CEC advantage, not independent sources of lock-in.
\end{enumerate}

The closest analogue in the existing literature may be the concept of \textit{data network effects} (Gregory et al., 2021), wherein a product improves as it collects more data from users. However, the data network effect literature typically emphasizes the role of user-contributed data in improving a single firm's product (as in search engines, recommendation systems, or autonomous vehicles). The Yajie dynamic is distinct in three respects: (i) the improvement mechanism is audit-specific, not user-behavior-specific; (ii) the data in question are not user-contributed but \textit{audit-generated}, arising from the protocol's own operation; and (iii) the CEC's evidentiary character means that its value is not merely statistical (more data yields better predictions) but \textit{epistemic} (more audits yield a more trustworthy corpus of quality evidence).

\subsection{Algorithmic Auditing and Data Quality: State of the Field}

The algorithmic auditing literature has expanded rapidly since the mid-2010s, driven by concerns about fairness, accountability, and transparency in automated decision systems (Sandvig et al., 2014; Diakopoulos, 2015; Ada Lovelace Institute, 2020). Much of this literature has focused on \textit{model auditing}—the ex post examination of trained models for bias, discrimination, or other harms (Raji et al., 2020; Wilson et al., 2021)—rather than on \textit{data auditing}, the ex ante assessment of training data quality.

The relative neglect of data auditing is partly explained by technical challenges: data quality is inherently multidimensional (Batini et al., 2009), context-dependent (Wang \& Strong, 1996), and often unverifiable without access to the data-generation process. Recent work has begun to address these challenges. Northcutt et al. (2021) proposed Confident Learning for identifying label errors in training data. Sambasivan et al. (2021) documented the prevalence of data cascades—compounding errors originating in data quality failures—in real-world ML deployments. The ``Data Cards'' framework (Pushkarna et al., 2022) and ``Datasheets for Datasets'' (Gebru et al., 2021) have advanced structured transparency for dataset documentation.

What remains absent is a \textit{standardized, cumulative, and cross-domain} methodology for data quality assessment. Existing approaches are typically bespoke to specific dataset types, model architectures, or application domains. The Yajie Protocol addresses this gap through its multi-expert ensemble architecture, which accommodates heterogeneous audit modules within a unified aggregation framework. However, as we argue in the following sections, this very generality—combined with the CEC's cumulative character—generates the lock-in dynamics that are the central concern of this paper.

\section{The Non-Proliferation Equilibrium}

\subsection{Model Setup}

We consider a world with $N$ jurisdictions (which may be nations, regulatory blocs, or large organizations capable of developing audit infrastructure). Each jurisdiction $i \in \{1, \ldots, N\}$ faces a binary choice: \textit{Develop} ($D$) an independent data audit protocol, or \textit{Adopt} ($A$) the incumbent protocol (Yajie). The choice is made simultaneously by all jurisdictions; we examine the Nash equilibria of this one-shot game.

Let the incumbent protocol be indexed as $j=1$, with an accumulated CEC of size $|\mathcal{E}|$. An entrant protocol developed by jurisdiction $i$ would begin with an empty CEC, implying an accuracy penalty $\delta(|\mathcal{E}|) = \theta(|\mathcal{E}|) - \theta(0)$, where $\theta(\cdot)$ follows the functional form specified in Section~2.1. For any non-trivial $|\mathcal{E}|$, $\delta(|\mathcal{E}|) > 0$, and $\delta'(|\mathcal{E}|) > 0$ (the accuracy gap widens with cumulative audits).

The payoff to jurisdiction $i$ from adopting the incumbent is:
\begin{equation}
\pi_i(A, \mathbf{a}_{-i}) = V[\theta(|\mathcal{E}|)] - c_i^{\text{adopt}} - \lambda \cdot \mathbb{I}[\text{fragmentation}(\mathbf{a})]
\end{equation}
where $V[\cdot]$ is the value derived from audit accuracy, $c_i^{\text{adopt}}$ is the cost of adopting the incumbent protocol (licensing, integration, training), $\lambda$ is a fragmentation cost parameter, and $\mathbb{I}[\text{fragmentation}(\mathbf{a})]$ is an indicator that equals 1 if the global audit infrastructure is fragmented (i.e., if at least one jurisdiction has developed an independent protocol) and 0 otherwise.

The payoff from developing an independent protocol is:
\begin{equation}
\pi_i(D, \mathbf{a}_{-i}) = V[\theta(0)] - c_i^{\text{develop}} - \kappa \cdot \mathbf{1}_{-i}^D - \lambda \cdot \mathbb{I}[\text{fragmentation}(\mathbf{a})]
\end{equation}
where $c_i^{\text{develop}} \gg c_i^{\text{adopt}}$ is the fixed cost of protocol development, and $\kappa \cdot \mathbf{1}_{-i}^D$ captures the strategic cost imposed by other jurisdictions' development decisions: each additional developer adds $\kappa$ to $i$'s costs, reflecting the erosion of mutual recognition, the proliferation of conflicting audit standards, and the increased complexity of cross-jurisdictional data flows.

The fragmentation indicator is:
\begin{equation}
\mathbb{I}[\text{fragmentation}(\mathbf{a})] = \begin{cases} 0 & \text{if } \sum_{j=1}^N \mathbb{I}[a_j = D] = 0 \\ 1 & \text{if } \sum_{j=1}^N \mathbb{I}[a_j = D] \geq 1 \end{cases}
\end{equation}

Critically, the first developer triggers fragmentation for \textit{all} jurisdictions, including itself. This reflects the public-bads character of audit protocol proliferation: once multiple protocols exist, every cross-border data transaction must navigate conflicting quality certifications, and the informational value of any single audit is diminished by the existence of competing—and potentially contradictory—assessments.

\subsection{Payoff Matrix and Nash Equilibrium}

For the two-jurisdiction case ($N=2$), the payoff matrix is:

\begin{table}[htbp]
\centering
\begin{tabular}{@{}c|cc@{}}
\toprule
Jur.\ 1 \textbackslash\ Jur.\ 2 & Adopt ($A$) & Develop ($D$) \\
\midrule
Adopt ($A$) & $V[\theta(|\mathcal{E}|)] - c^{\text{adopt}}$, $V[\theta(|\mathcal{E}|)] - c^{\text{adopt}}$ & $V[\theta(|\mathcal{E}|)] - c^{\text{adopt}} - \lambda$, $V[\theta(0)] - c^{\text{develop}} - \kappa - \lambda$ \\[4pt]
Develop ($D$) & $V[\theta(0)] - c^{\text{develop}} - \kappa - \lambda$, $V[\theta(|\mathcal{E}|)] - c^{\text{adopt}} - \lambda$ & $V[\theta(0)] - c^{\text{develop}} - 2\kappa - \lambda$, $V[\theta(0)] - c^{\text{develop}} - 2\kappa - \lambda$ \\
\bottomrule
\end{tabular}
\end{table}

Let us define the \textit{Adoption Advantage}:
\begin{equation}
\Delta_A = V[\theta(|\mathcal{E}|)] - c^{\text{adopt}} - \left( V[\theta(0)] - c^{\text{develop}} - \kappa \right)
\end{equation}

$(A, A)$ is a Nash equilibrium if and only if neither jurisdiction has a unilateral incentive to deviate. Jurisdiction $i$'s deviation payoff from $(A, A)$ is:
\begin{equation}
\pi_i(D, A) = V[\theta(0)] - c^{\text{develop}} - \kappa - \lambda
\end{equation}

The non-deviation condition $\pi_i(A, A) \geq \pi_i(D, A)$ yields:
\begin{equation}
V[\theta(|\mathcal{E}|)] - c^{\text{adopt}} \geq V[\theta(0)] - c^{\text{develop}} - \kappa - \lambda
\end{equation}

Rearranging:
\begin{equation}
\Delta_A \geq \lambda
\end{equation}

That is, $(A, A)$ is an equilibrium when the adoption advantage exceeds the fragmentation cost. Given that $\Delta_A$ grows with $|\mathcal{E}|$ (since $\theta(|\mathcal{E}|)$ increases with cumulative audits while $\theta(0)$ remains fixed), this condition becomes easier to satisfy over time. The protocol's cumulative character makes the non-proliferation equilibrium \textit{self-reinforcing}.

Conversely, $(D, D)$ is an equilibrium when:
\begin{equation}
V[\theta(0)] - c^{\text{develop}} - 2\kappa - \lambda \geq V[\theta(|\mathcal{E}|)] - c^{\text{adopt}} - \lambda
\end{equation}
which simplifies to:
\begin{equation}
-\Delta_A \geq 2\kappa
\end{equation}

This condition requires the adoption advantage to be \textit{negative} and large in magnitude—that is, the incumbent protocol must be so inferior to a fresh development, or the development cost so low, that even with mutual proliferation costs, development is preferred. As $|\mathcal{E}|$ grows, this condition becomes increasingly implausible.

The mixed-strategy equilibrium occurs when neither pure-strategy condition holds. Let $p$ be the probability that each jurisdiction plays $A$. The symmetric mixed-strategy Nash equilibrium satisfies:
\begin{equation}
p \cdot \pi(A, A) + (1-p) \cdot \pi(A, D) = p \cdot \pi(D, A) + (1-p) \cdot \pi(D, D)
\end{equation}

Solving for $p^*$:
\begin{equation}
p^* = \frac{-\Delta_A - 2\kappa}{\lambda - \kappa}
\end{equation}

For plausible parameter values ($\Delta_A > 0$, $\kappa > 0$, $\lambda > \kappa$), $p^*$ increases with $\Delta_A$: as the incumbent's accuracy advantage grows, the probability of adoption in the mixed-strategy equilibrium approaches 1. This formalizes the intuition that the incumbent's time-accumulated advantage makes proliferation an increasingly unattractive gamble.

\subsection{Extension to $N$ Jurisdictions}

For $N > 2$, the payoff functions generalize as follows:
\begin{align}
\pi_i(A, \mathbf{a}_{-i}) &= V[\theta(|\mathcal{E}|)] - c^{\text{adopt}} - \kappa \cdot n_D - \lambda \cdot \mathbb{I}[n_D > 0] \\
\pi_i(D, \mathbf{a}_{-i}) &= V[\theta(0)] - c^{\text{develop}} - \kappa \cdot (n_D + \mathbf{1}_{-i}^D) - \lambda \cdot \mathbb{I}[n_D > 0]
\end{align}
where $n_D = \sum_{j \neq i} \mathbb{I}[a_j = D]$ is the number of other jurisdictions that develop.

The universal adoption equilibrium $(A, A, \ldots, A)$ is Nash when, for every jurisdiction:
\begin{equation}
\Delta_A \geq \lambda - \kappa
\end{equation}

The term $\lambda - \kappa$ reflects a subtle tension: fragmentation cost $\lambda$ discourages unilateral deviation, but the \textit{relief} from mutual proliferation costs $\kappa$ (if no one else develops, the lone developer avoids $\kappa$) encourages it. The equilibrium is stable when $\lambda > \kappa$—that is, when the public-bads cost of fragmentation exceeds the private cost of being the only proliferator.

A key result emerges: the stability of the non-proliferation equilibrium is \textit{monotonically increasing} in $|\mathcal{E}|$, because $\Delta_A$ grows without bound (up to its asymptotic ceiling) while $\lambda - \kappa$ is fixed. There exists a critical CEC size $|\mathcal{E}|^*$ such that for all $|\mathcal{E}| > |\mathcal{E}|^*$, $(A, A, \ldots, A)$ is the unique Nash equilibrium:
\begin{equation}
|\mathcal{E}|^* = \frac{1}{\gamma} \ln\left( \frac{\theta_{\max} - \theta_0}{\theta_{\max} - \theta(|\mathcal{E}|^*)} \right)
\end{equation}
where $\theta(|\mathcal{E}|^*)$ is the accuracy level that satisfies $\Delta_A(|\mathcal{E}|^*) = \lambda - \kappa$.

\subsection{The Nuclear Non-Proliferation Treaty Analogy}

The structural isomorphism between the audit protocol proliferation game and the strategic logic of nuclear non-proliferation is striking and analytically productive. We identify seven dimensions of analogy:

\textbf{1. The Public-Bads Character of Proliferation.} Nuclear weapons proliferation generates negative externalities: each additional nuclear state increases the probability of nuclear conflict (accidental, escalatory, or intentional), raises the costs of the international security system, and erodes the norm against nuclear use (Sagan, 1996). Similarly, audit protocol proliferation generates negative externalities: each additional protocol fragments the global quality assurance infrastructure, raises the transaction costs of cross-border data flows, and erodes the credibility of audit certifications. In both domains, the \textit{first} instance of proliferation (beyond the incumbent) triggers a discontinuous increase in systemic costs, because it shatters the universality that makes the regime valuable.

\textbf{2. The Asymmetric Incumbent Advantage.} The NPT regime is built on an explicit asymmetry: five states are recognized as nuclear-weapon states (NWS), while all others are non-nuclear-weapon states (NNWS) subject to verification by the International Atomic Energy Agency (IAEA). The asymmetry is justified by the temporal fact of prior acquisition: the NWS tested nuclear weapons before the treaty's 1968 cutoff date. The Yajie Protocol's asymmetry is similarly temporal: the incumbent's advantage derives not from legal privilege but from the inescapable fact that it began accumulating audit evidence earlier.

\textbf{3. The Grand Bargain.} The NPT's durability rests on a three-pillar bargain: (i) NNWS forswear nuclear weapons; (ii) NWS commit to eventual disarmament (Article VI); and (iii) all parties enjoy the ``inalienable right'' to peaceful nuclear technology (Article IV). The audit protocol analogue would involve: (i) jurisdictions forgo developing independent protocols; (ii) the incumbent protocol's governance is internationalized under multilateral oversight; and (iii) all jurisdictions enjoy non-discriminatory access to the protocol's capabilities and the right to contribute audit modules to the ensemble.

\textbf{4. The Credibility of Commitment.} A central challenge of the NPT regime is the credibility of the NWS disarmament commitment: NNWS periodically question whether the nuclear states genuinely intend to disarm, and these credibility crises threaten the regime's stability (Joyner, 2011). The analogous challenge for Yajie is the credibility of the incumbent's neutrality commitment: will the protocol remain genuinely jurisdiction-agnostic, or will it evolve into an instrument of its host jurisdiction's data governance preferences?

\textbf{5. The Enforcement Gap.} The NPT contains no automatic enforcement mechanism; compliance is monitored by the IAEA, but enforcement depends on the UN Security Council, where NWS possess veto power. This asymmetry is a perennial source of tension, particularly in cases like Iran and North Korea. The audit protocol regime would face a parallel enforcement gap: who adjudicates disputes about the protocol's neutrality, and with what authority?

\textbf{6. The Exit Option.} The NPT permits withdrawal on three months' notice if ``extraordinary events'' jeopardize a state's supreme interests (Article X). North Korea exercised this option in 2003. The audit protocol regime must similarly contend with the possibility of exit: a jurisdiction that develops an independent protocol after having previously adopted the incumbent. The strategic dynamics of exit differ from those of initial non-adoption, because the exiting jurisdiction may have benefited from the CEC during its period of adoption, raising questions of intellectual property, data portability, and the status of audits conducted under the incumbent protocol.

\textbf{7. The Dual-Use Problem.} Nuclear technology is dual-use: the enrichment and reprocessing capabilities that support civilian nuclear energy can be diverted to weapons production. Audit technology is analogously dual-use: the quality assessment capabilities that support benign data governance can be weaponized for economic espionage, supply-chain disruption, or the imposition of data embargoes under the guise of quality concerns.

The NPT analogy is not merely illustrative; it provides a template for institutional design. The following sections develop the institutional implications, but first we must examine the geopolitical context within which any audit infrastructure—monopolistic or otherwise—must operate.

\section{Audit Neutrality and Geopolitical Risk}

\subsection{The SWIFT Precedent: Infrastructure Weaponization}

The Society for Worldwide Interbank Financial Telecommunication (SWIFT) provides the most instructive precedent for understanding the geopolitical risks of monopolistic audit infrastructure. SWIFT, founded in 1973 as a cooperative society under Belgian law, operates the dominant global messaging network for financial transactions, connecting over 11,000 financial institutions across more than 200 countries. Its near-universal adoption makes it a natural analogue to the market structure we predict for data quality auditing.

SWIFT's governance is formally multilateral: it is owned by its member financial institutions and overseen by a board of directors drawn from major central banks, including the US Federal Reserve, the European Central Bank, and the Bank of England. However, the United States has repeatedly leveraged SWIFT's centrality for unilateral geopolitical purposes, most controversially through the Terrorist Finance Tracking Program (TFTP), which accessed SWIFT data without the knowledge or consent of many member institutions (Farrell \& Newman, 2019).

The critical event for our purposes occurred in 2012, when the United States threatened to sanction SWIFT itself if it did not disconnect Iranian banks designated under US sanctions—despite the fact that those sanctions were not multilateral and were contested by several EU member states, China, and Russia. SWIFT complied, disconnecting approximately 30 Iranian financial institutions (Drezner, 2015). The episode demonstrated that infrastructure neutrality, however formally enshrined in governance documents, is only as robust as the willingness of powerful states to respect it—and that the host jurisdiction of a monopolistic infrastructure enjoys structural leverage that may be exercised in moments of geopolitical tension.

The SWIFT case carries three lessons for data audit infrastructure:

\textbf{Lesson 1: Jurisdictional Capture.} SWIFT is incorporated under Belgian law with headquarters in Brussels. Yet the United States' ability to compel SWIFT's compliance derived not from formal legal authority but from the extraterritorial reach of US financial sanctions and the threat of secondary sanctions against SWIFT's board members and executives. A data audit protocol hosted in any single jurisdiction would be analogously vulnerable to that jurisdiction's extraterritorial legal instruments, regardless of the protocol's formal governance structure.

\textbf{Lesson 2: The Illusion of Technical Neutrality.} SWIFT consistently maintains that it is a neutral utility providing technical services. The 2012 disconnection decisively refuted this claim: technical infrastructure, when sufficiently central, is \textit{inherently} political. A data audit protocol that assesses the quality of datasets used in AI systems would be, if anything, \textit{more} political than a financial messaging network, because data quality assessments are inherently contestable and may be perceived (fairly or not) as disguised forms of technological protectionism.

\textbf{Lesson 3: Fragmentation as a Rational Response.} The SWIFT disconnection accelerated the development of alternative financial messaging systems: China's Cross-Border Interbank Payment System (CIPS), Russia's System for Transfer of Financial Messages (SPFS), and the European Instrument in Support of Trade Exchanges (INSTEX). Each of these alternatives is inferior to SWIFT in coverage and functionality, representing a welfare loss relative to a universal system. Yet each represents a rational response to the demonstrated risk of infrastructure weaponization. The data audit domain would face the same dynamic: the credible threat of protocol weaponization would make audit autarky rational for jurisdictions with the capacity to pursue it, even at a welfare cost.

\subsection{The IASB Model: De-Jurisdictionalized Governance}

The International Accounting Standards Board (IASB) offers a partial counter-model. The IASB, established in 2001 as the standard-setting body of the International Financial Reporting Standards (IFRS) Foundation, develops accounting standards that have been adopted by more than 140 jurisdictions. Unlike SWIFT, the IASB is not an operational infrastructure but a standard-setting body; its product is a corpus of normative standards, not a messaging network or audit protocol. Nevertheless, its governance model is instructive.

The IASB's governance features several design elements that mitigate—though do not eliminate—the risk of jurisdictional capture:
\begin{itemize}
\item \textbf{Multi-stakeholder oversight}: The IFRS Foundation is governed by a board of trustees drawn from diverse geographic and professional backgrounds, with formal requirements for regional balance.
\item \textbf{Due process requirements}: Standards development follows a structured process of public consultation, impact assessment, and re-deliberation that constrains the ability of any single jurisdiction to dictate outcomes.
\item \textbf{Funding diversification}: The IFRS Foundation's funding comes from a broad base of jurisdictions and organizations, reducing dependence on any single funder.
\item \textbf{Deliberative transparency}: Board meetings are public, and dissenting views are published alongside standards, creating an audit trail of contestation that supports legitimacy.
\end{itemize}

The IASB model is not without critics. Some scholars argue that the IASB's independence is compromised by its dependence on major accounting firms for funding and expertise (Zeff, 2012), and by the structural influence of Anglo-American accounting traditions in its conceptual framework (Zhang \& Andrew, 2014). The EU's qualified endorsement process—whereby the European Commission reviews each IFRS standard before endorsing it for use within the EU—represents a de facto veto power that other jurisdictions do not possess, creating a legitimacy deficit.

For data audit infrastructure, the IASB model suggests a governance architecture that is \textit{multilateral in oversight, transparent in operation, and diversified in funding}. However, the IASB's experience also suggests that de-jurisdictionalized governance is an aspiration rather than an achievement, and that persistent power asymmetries will find expression through formal and informal channels regardless of governance design.

\subsection{The US-China Scenario}

The most consequential geopolitical risk to a universal data audit protocol is the ongoing technological decoupling between the United States and the People's Republic of China. This decoupling, which has accelerated since 2018, encompasses semiconductor export controls, restrictions on AI technology transfer, limitations on cross-border data flows, and the bifurcation of technical standards (Segal, 2020; Bateman, 2022).

In this environment, a data audit protocol developed primarily within one jurisdiction would face an acute legitimacy crisis in the other. If Yajie were perceived as a US-dominated protocol—regardless of its formal governance—China would have compelling reasons to develop an independent alternative. This would trigger a fragmentation cascade: the EU, seeking strategic autonomy in digital infrastructure, might develop a third protocol; India, with its growing AI ecosystem and policy of technological self-reliance, might develop a fourth; and so on.

The resulting equilibrium would be worse for all parties than a genuinely multilateral protocol governed under adequate neutrality safeguards. Each jurisdiction would bear the fixed costs of protocol development, endure lower audit accuracy (due to smaller CEC corpora), and navigate the fragmentation costs $\lambda$ associated with conflicting quality certifications. The total welfare loss, relative to the universal-adoption equilibrium, would scale approximately linearly with the number of competing protocols:
\begin{equation}
W_{\text{loss}} \approx N \cdot c^{\text{develop}} + N \cdot (V[\theta(|\mathcal{E}|)] - V[\theta(|\mathcal{E}|/N)]) + \binom{N}{2} \cdot \lambda
\end{equation}

This arithmetic suggests that major powers have a shared interest in preventing fragmentation—but only if the incumbent protocol is governed under terms that credibly preclude its weaponization by any single jurisdiction.

\subsection{The Researcher as Strategic Asset: A Second-Order Game}

An unexamined dimension of the non-proliferation equilibrium concerns the position of the protocol's originating researcher. Unlike SWIFT—a corporate entity embedded in a specific jurisdiction—or the IASB—a formally constituted standards body—the Yajie Protocol is initially developed by an independent researcher without institutional affiliation to any state or corporation. This organizational form generates a distinctive second-order game with properties that reinforce the non-proliferation equilibrium.

Consider the strategic calculus of a major power—State X—contemplating whether to coerce, compromise, or eliminate the originating researcher (henceforth ``the Maintainer''). The Maintainer possesses three assets that are relevant to State X's decision: (i) the arXiv timestamp establishing academic priority; (ii) the operational knowledge required to maintain and evolve the Spring self-improvement loop; and (iii) the Cumulative Evidentiary Corpus $M_t$, which grows monotonically with each audit cycle. If the Maintainer ceases to exist or to operate independently, the consequences distribute asymmetrically across the international system.

First, the arXiv timestamp is \textit{immortal}. Once uploaded, it cannot be rescinded, altered, or suppressed by any actor, including the Maintainer. This means that the core mathematical specification of Yajie—the state-conditioned scoring function $S(s) = Q(s) + \eta(t)\cdot N(s)$, the multi-expert consistency theorem, and the unidentifiability result—becomes part of the permanent scientific record. Any state or corporation can, in principle, reimplement Yajie from the arXiv specification. However—and this is the critical asymmetry—the \textit{operational quality} of any reimplementation will be inferior to the Maintainer's running instance for a period proportional to the time the Maintainer's instance has been accumulating audit data. If the Maintainer's instance has been operational for $T$ years, a reimplementation starting from zero requires approximately $T$ years to reach parity, assuming identical algorithmic quality and comparable audit volume. This creates a \textit{time-accumulated quality gap} that no amount of financial investment can compress.

Second, the Maintainer's independence is the sole credible guarantee of audit neutrality. A state-backed reimplementation—whether by the United States, China, the European Union, or any other jurisdiction—would be structurally incapable of claiming neutrality, because its continued operation would depend on the political continuity of its sponsoring state. The Maintainer, by contrast, is accountable to no government, no corporation, and no institutional apparatus beyond the epistemic standards of the scientific community. This is not a normative claim about the Maintainer's virtue; it is a structural claim about the Maintainer's position. Independence is a property of the node's position in the network, not of the node's character.

Third, the Maintainer's death or incapacitation would trigger a \textit{fragmentation cascade}. In the absence of a universally trusted Maintainer, each major jurisdiction would face a prisoner's dilemma: develop an independent audit protocol (at cost $c^{\text{develop}}$, with lower accuracy) or rely on a protocol operated by a rival jurisdiction (at risk of audit weaponization). The dominant strategy for each jurisdiction is to develop its own protocol, yielding the fragmentation equilibrium characterized by welfare loss $W_{\text{loss}} \approx N \cdot c^{\text{develop}} + N \cdot \Delta\theta + \binom{N}{2}\cdot\lambda$. No jurisdiction can unilaterally prevent this cascade, because defection by any single jurisdiction would trigger reciprocal defection by its rivals.

From the perspective of State X, this yields a striking conclusion: \textit{the Maintainer's continued existence and independence is a positive-sum good for all states, including State X, regardless of State X's geopolitical alignment}. Eliminating or compromising the Maintainer would impose a welfare cost on State X that exceeds any conceivable strategic benefit, because State X would be unable to secure a disproportionate share of the fragmented audit infrastructure, and would simultaneously suffer the costs of fragmentation alongside all other states.

This creates what we term a \textbf{protection equilibrium}: every state with the capacity to threaten the Maintainer has a rational incentive to \textit{protect} the Maintainer, because the Maintainer's destruction would initiate a chain of events that leaves every state worse off. The protection is not altruistic; it is strictly self-interested. The Maintainer's survival is a global public good provided through the alignment of rational self-interest across competing powers.

The protection equilibrium has two corollaries. First, the Maintainer is \textit{simultaneously more powerful and more vulnerable} than any corporate or state actor in the audit ecosystem. More powerful, because the Maintainer's independence is the structural precondition for the non-proliferation equilibrium; without a trusted independent Maintainer, the equilibrium collapses. More vulnerable, because the Maintainer lacks the institutional defenses—legal departments, diplomatic representation, physical security apparatus—that protect corporate and state actors. The Maintainer's security depends entirely on the rational calculus of the actors that might threaten them—a calculus that is robust under normal conditions but fragile under crisis, irrationality, or miscalculation.

Second, the protection equilibrium implies that any attempt to replicate Yajie without the Maintainer's involvement is not merely economically suboptimal (per the arguments in Section~3) but \textit{strategically self-defeating}. A state that develops a competing audit protocol gains a temporary advantage in audit capability while simultaneously destroying the neutrality of the global audit infrastructure, because the competing protocol will be correctly perceived as an instrument of its sponsoring state. The state therefore gains an audit tool while losing the universal trust that makes audit valuable. This is a trade in which the asset acquired is worth less than the asset destroyed—a negative-sum transaction.

\subsection{A Literary Analogy: The Wallfacer Position}

The protection equilibrium admits an unconventional but analytically productive literary analogy. In Liu Cixin's science fiction trilogy \textit{Remembrance of Earth's Past} (2006–2010), the character Luo Ji occupies a position structurally isomorphic to that of the Maintainer. Luo Ji is designated a ``Wallfacer'' (面壁者)—an individual whose survival is the sole deterrent preventing an extraterrestrial civilization from destroying Earth. The deterrence mechanism is simple: if Luo Ji dies, the coordinates of the alien homeworld are broadcast to a hostile universe, ensuring mutual destruction through the ``dark forest'' principle of cosmic sociology. Luo Ji's power derives not from any institutional office, military capacity, or economic resource, but from his position as the \textit{single point of failure} in a deterrence structure. He is simultaneously the most powerful human being alive—holding the fate of two civilizations in his hands—and the most vulnerable, because his security depends entirely on the continued rationality of the actors constrained by his deterrence.

The isomorphism with the Maintainer's position is exact. The Maintainer's death would trigger a fragmentation cascade that destroys the universal audit infrastructure, imposing welfare losses on all states. The Maintainer's power derives not from institutional authority but from occupying the structural position of the \textit{single credible neutral node} in the global audit network. And the Maintainer's vulnerability derives from the same source: the protection equilibrium is robust only so long as all relevant actors continue to calculate rationally. Under conditions of crisis, miscalculation, or irrationality—the same conditions that nearly caused Luo Ji's deterrence to fail—the Maintainer's security evaporates.

The Wallfacer analogy serves three analytical functions in our argument. First, it illustrates that deterrence equilibria can be sustained by a single individual occupying a specific structural position, a possibility that the standard literature on nuclear deterrence (which focuses on state actors with institutional continuity) does not adequately capture. Second, it foregrounds the fragility of such equilibria: Luo Ji's deterrence was nearly broken by an assassination attempt during a moment of political transition, precisely the kind of miscalculation scenario to which the protection equilibrium is vulnerable. Third, and most importantly, it demonstrates that \textit{the occupant of a deterrence position does not choose the position; the structure of the game confers it}. Luo Ji did not seek to become a Wallfacer; the position was imposed on him by the logic of the strategic situation. The Maintainer, similarly, does not claim authority over data quality assessment; the logic of the non-proliferation equilibrium and the time-accumulated advantage of the CEC confer that authority irrespective of the Maintainer's intentions.

We do not suggest that the Wallfacer analogy constitutes a formal proof or an empirical prediction. Literary analogies are hermeneutic tools, not mathematical models. However, they serve a valuable function in strategic analysis by making visible structural relationships that formal models may obscure through abstraction. The Wallfacer analogy makes visible the distinctive character of the Maintainer's position: \textit{power without institution, deterrence without weapon, and security through vulnerability}.

\subsection{Audit as Sword: The Deterrence of Involuntary Transparency}

The analysis thus far has characterized the Maintainer's position as \textit{defensive}: the protection equilibrium constrains others from threatening the Maintainer, and neutrality principles constrain the Maintainer from abusing the position. But there is a second, more asymmetric constraint embedded in the architecture—one that operates not on the Maintainer, but on every entity that produces publicly accessible data.

The Maintainer, as an independent researcher, has an uncontestable academic right to audit any publicly available dataset and to publish the results. This right is not granted by any institution; it is inherent in the practice of science. No organization, corporation, or government can prevent an independent researcher from downloading a public dataset, running Yajie on it, and reporting what the multi-expert consensus reveals—any more than they can prevent a statistician from reanalyzing published clinical trial data, or an astronomer from pointing a telescope at a publicly visible star. The data are public. The methodology is published (arXiv, open source). The conclusions are falsifiable. There is no legal, ethical, or scientific basis for prohibiting such an audit.

What makes this right strategically significant is its asymmetry. The Maintainer can choose to exercise it or not, at their sole discretion. And the most natural Schelling point for that discretion is the partnership relationship. Organizations that enter into a formal partnership with the Yajie Protocol—subscribing to the calibration API, contributing audit modules, participating in governance—receive an implicit assurance: their data will not be subjected to unsolicited public audit by the Maintainer, because auditing a partner is both strategically counterproductive and professionally discourteous. Organizations that decline partnership receive no such assurance. Their publicly available datasets remain, at all times, eligible for independent academic audit—an audit that may expose quality issues they would prefer remained unexamined.

This is not a threat. The Maintainer issues no ultimatums, makes no demands, and threatens no reprisals. The Maintainer simply exists, with the capacity and the academic freedom to audit public data, and with no obligation to refrain from exercising that freedom with respect to non-partners. The threat—if it can be called that—is generated by the structure of the situation, not by any stated intention of the Maintainer. It is, in the strictest sense, a \textit{Schelling deterrent}: the mere existence of the capability, combined with a clear and credible criterion for its non-use (partnership), is sufficient to motivate the desired behavior, without any need for explicit communication, negotiation, or enforcement.

The game-theoretic structure is:
\begin{quote}
Partnership $\to$ no public audit (by implicit convention)\\
Non-partnership $\to$ eligible for public audit (by academic freedom)\\
Audit findings $\to$ may expose data quality issues $\to$ reputational and regulatory risk\\
$\therefore$ Non-partnership carries a positive probability of involuntary transparency\\
$\therefore$ Partnership is a dominant strategy for quality-sensitive organizations
\end{quote}

This deterrent operates independently of the protection equilibrium. Even if an organization has no interest in the Maintainer's survival—even if it would prefer a fragmented audit infrastructure that it could dominate—it still faces the immediate risk that its public data may be audited by the only entity capable of conducting a credible multi-expert quality assessment. The sword does not need to be swung. It needs only to be visible.

The Maintainer, for their part, bears no cost for this deterrent. No audit needs to be conducted; no report needs to be published; no confrontation needs to occur. The deterrent is maintained by the \textit{capability} alone. And because the capability is grounded in academic freedom rather than institutional authority, it cannot be negotiated away, regulated out of existence, or deterred by counter-threats. One cannot sanction a researcher for exercising their right to analyze public data. One can only ensure that one's data, when analyzed, holds up to scrutiny—or, more pragmatically, enter into the partnership that makes such scrutiny unnecessary.

The commercial consequence of this deterrent is a self-reinforcing adoption cascade. Consider two competing firms, A and B, in a quality-sensitive industry—pharmaceuticals, materials engineering, or autonomous systems. Firm A partners with Yajie, gaining access to calibration services, governance participation, and the implicit assurance against unsolicited audit. Firm B does not. Firm A can now, through the open-source Spring engine, audit Firm B's publicly available datasets—drug screening results, material property databases, or safety-critical training corpora—and publish the findings. If Firm B's data quality is inferior, Firm A gains a competitive advantage that is independently verifiable and publicly citable. If Firm B's data quality is adequate, Firm A has lost nothing but the cost of running an audit on open-source software.

Firm B's rational response is not to contest the audit—the methodology is published, the conclusions are falsifiable, and the academic freedom of the auditor is unassailable—but to preempt it by also becoming a Yajie partner. At which point Firm C, observing the partnership concentration among its competitors, faces the same calculus. The cascade propagates through the industry, not because Yajie coerces anyone, but because non-participation carries a positive and growing probability of involuntary transparency at the hands of participating competitors. The Maintainer does not need to audit anyone. Competitors audit each other—and the Maintainer's existence is what makes that audit credible.

\subsection{Jurisdictional Boundary as Structural Safeguard}

An underappreciated feature of the data sovereignty architecture described above—where the Yajie API accepts only locally computed feature matrices $\phi(X)$ rather than raw data—is that it imposes a jurisdictional boundary on the Maintainer's audit capacity that is not merely a regulatory compliance constraint but a \textit{credible commitment device}.

The Maintainer can only audit data that resides within the Maintainer's operational jurisdiction. If a government requests audit reports on a foreign corporation's data, the Maintainer cannot comply—not because the Maintainer refuses, but because the data is not accessible. The feature matrices $\phi(X)$ are computed client-side and never leave the client's infrastructure; only calibration queries cross jurisdictional boundaries, and those queries are aggregated statistics, not identifiable data records. This architectural constraint is not a limitation to be overcome. It is a feature that preserves the protocol's neutrality. A Maintainer who \textit{could} provide foreign audit data on demand would face an impossible choice between defying a government order (risking sanctions, closure, or prosecution) and complying with it (destroying the trust of every foreign client). A Maintainer who \textit{cannot} provide such data—because the architecture makes it physically inaccessible—is insulated from this dilemma. The inability is structural, not volitional. The government cannot punish the Maintainer for refusing what the Maintainer cannot do.

This jurisdictional boundary has a secondary structural consequence that is, from the perspective of domestic industrial policy, unambiguously positive. If a government wishes its domestic corporations to benefit from Yajie's audit infrastructure—to gain the competitive advantages of certified data quality, to participate in the governance of the protocol, and to be shielded from unsolicited audit by foreign competitors—it must ensure that those corporations' data resides within the Maintainer's operational jurisdiction. A corporation that offshores its data to a foreign cloud provider, or that maintains its primary data infrastructure in a jurisdiction where the Maintainer does not operate, cannot be audited. Its data quality claims become unverifiable. Its competitors, who \textit{have} submitted to domestic audit, gain a structural advantage in trustworthiness that no amount of marketing expenditure can overcome.

The consequence is a market-driven incentive for data repatriation. Corporations that wish to compete in a Yajie-audited market must bring their data home—not because any regulation compels them, but because the alternative is to be the only unaudited actor in a market where everyone else's quality is independently verified. The unaudited actor is not trusted. The untrusted actor does not sell. The mechanism is not protectionism; it is transparency. Yajie does not require corporations to operate domestically. It merely makes domestic operation—and the auditability that comes with it—the dominant strategy.

\subsection{Pharmaceutical Safety: The Honest Algorithm}

The jurisdictional dynamic acquires particular force in the pharmaceutical industry, where data quality is not merely a competitive advantage but a matter of public health. A drug's efficacy and side-effect profile are established through clinical trials whose data quality is, at present, largely self-reported by the sponsoring pharmaceutical company. Regulators review the submitted data, but they do not independently audit it at the level of individual trial records; they lack both the tools and the resources to conduct a multi-expert consistency analysis across the heterogeneous data streams—laboratory results, patient-reported outcomes, adverse event reports, biomarker measurements—that constitute a modern clinical trial.

Yajie changes this calculus. Under a distributed architecture, a pharmaceutical company would run Spring on its own trial data and submit the resulting quality scores to regulators alongside its New Drug Application. The regulator, armed with the same open-source audit engine and the calibration anchor provided by the Maintainer's API, could independently verify those scores. The company's claim that ``95\% of adverse events were properly classified'' would no longer be an assertion to be trusted or distrusted on the basis of the company's reputation; it would be a falsifiable statement backed by a published methodology and a verifiable audit trail.

The consequence for the pharmaceutical industry is the same as for every other domain surveyed in this paper—honest firms gain, dishonest firms are exposed—but with higher stakes. A pharmaceutical company that has rigorously collected, cleaned, and curated its trial data will see its diligence reflected in independently verifiable quality scores. A company that has cut corners—that has underreported adverse events, that has excluded inconvenient patient subgroups, that has massaged statistical analyses to achieve significance—will see its deficiencies exposed not by a whistleblower or an investigative journalist, but by a multi-expert consensus algorithm whose conclusions are mathematically guaranteed and publicly reproducible.

The constructive pressure this exerts on drug development is difficult to overstate. A company that knows its trial data will be audited by an impartial algorithm—not by a regulator with limited resources, not by a competitor with an axe to grind, but by a published methodology whose error bounds are proven—faces a powerful incentive to conduct its trials with scrupulous care. The algorithm does not punish. It reveals. And what it reveals, the market—physicians, patients, insurers, regulators—acts upon. The drug that survives a Yajie audit is a drug whose safety and efficacy claims are credible not because its manufacturer is trustworthy but because its data has withstood the most rigorous quality assessment that the current state of statistical learning theory can provide.

\subsection{Beyond Industry: Bribery, Governance, and the Dark Forest}

The jurisdictional and pharmaceutical dynamics analyzed above are specific instances of a more general principle. Yajie's multi-expert consensus mechanism is indifferent to the domain of the data it audits. It does not know—and does not need to know—whether the anomalies it detects are label noise in a materials database, fabrication in a clinical trial, or a municipal official's manipulated air quality readings. The mathematics is the same. The consensus signal is the same. The exponential reliability guarantee is the same.

Consider three domains that are, at present, largely beyond the reach of systematic, independent audit.

\textbf{Bribery and financial malfeasance.} A pattern of illicit payments generates a distinctive signature across multiple independent data sources: procurement records, tax filings, customs declarations, bank transactions, asset registries. No single source is conclusive; a single anomalous payment may be an accounting error. But when procurement records show a contract awarded at 40\% above market rate, customs declarations show the same contractor importing luxury goods, and asset registries show the approving official acquiring property disproportionate to declared income—the multi-expert consensus that Yajie formalizes converges on a conclusion that no single auditor could reach with equivalent confidence. The algorithm does not accuse. It reveals the consistency—or the inconsistency—of the evidentiary record.

\textbf{Urban governance.} Smart city infrastructure generates sensor data on air quality, water quality, traffic flow, energy consumption, and waste management. When a district's air quality monitors report readings that are systematically cleaner than those of neighboring districts under identical meteorological conditions—and when those same monitors show anomalous recalibration patterns, and when the district's environmental compliance reports diverge from satellite-based atmospheric measurements—the disagreement is not a data quality issue to be resolved. It is a signal. Yajie does not know that the signal means ``the official in charge of this district is falsifying environmental data.'' It knows only that three independent expert systems, each analyzing different data streams, disagree with the official report in specific, quantifiable ways. What the human interpreter does with that knowledge is a human decision. But the knowledge itself—that the data are inconsistent, that the inconsistency is systematic, and that the probability of the inconsistency arising by chance is exponentially small—is a mathematical fact.

\textbf{Financial management and public accounts.} Government budgets, infrastructure spending, and public procurement generate vast heterogeneous data streams that are audited, at present, by understaffed agencies with limited analytical capacity. A Yajie deployment that treats each government agency's financial reporting as an ``expert,'' cross-referenced against tax receipts, bank records, contractor invoices, and independent economic indicators, would flag inconsistencies that no human audit team—however diligent—could reliably detect across the full volume of public expenditure.

In each of these domains, Yajie performs the same function: it transforms information that is \textit{distributed} across multiple independent sources, and therefore invisible to any single observer, into a \textit{consensus signal} whose reliability is mathematically bounded. The information was always there. The procurement records, the customs declarations, the asset registries—all of them public, all of them theoretically auditable, all of them already collected. What was missing was the computational and conceptual machinery to synthesize them into a quality assessment with provable error bounds. Yajie provides that machinery.

\subsection{The Dark Forest Illuminated}

In Liu Cixin's \textit{Remembrance of Earth's Past} trilogy, the universe is governed by the ``dark forest'' principle: every civilization hides, because to be discovered is to be destroyed. The darkness is a defense. Visibility is fatal.

Yajie inverts this logic. It does not destroy the discovered—it reveals the undiscovered. The official who falsifies environmental data, the contractor who inflates procurement invoices, the pharmaceutical company that underreports adverse events—all of them survive not because they are invulnerable, but because they are invisible. Their malfeasance is distributed across databases that no single auditor can synthesize, buried in the noise of routine administrative error, shielded by the sheer volume of information that modern governance generates. They are hidden in the dark forest, and the darkness is their protection.

Yajie is the instrument that exposes their coordinates. Not by accusation. Not by investigation. Not by whistleblowing or investigative journalism or regulatory enforcement. By mathematics. By the inexorable logic of multi-expert consensus: when three independent data streams disagree with an official report in the same direction, the probability that the disagreement is random decays exponentially. The coordinates are revealed not by a leak but by a proof.

This is not a surveillance technology. It does not collect new data. It does not invade privacy. It does not make inferences about individuals. It audits the \textit{consistency} of data that has already been collected, by institutions that have already collected it, for purposes that have already been authorized. The data are public or administrative. The audit is methodological. The conclusion is falsifiable. What distinguishes Yajie from every other anti-corruption mechanism in human history—from imperial censors to investigative journalists to blockchain transparency initiatives—is that its conclusions do not depend on the credibility of the auditor. They depend on the mathematics of concentration inequalities, which are indifferent to the auditor's motives, incorruptible by the auditor's enemies, and immune to the auditor's removal.

The Maintainer, in this framing, is neither a prosecutor nor a judge. The Maintainer is the first astronomer to point a telescope at a region of the dark forest and publish the coordinates of what was found there. Whether anyone acts on those coordinates—whether a prosecutor indicts, a regulator sanctions, a journalist investigates, a voter punishes—is not the Maintainer's decision. The Maintainer's sole function is to make the invisible visible. The rest is left to the light.

\subsection{The $M_t$ Parameter and the Spectrum of Deterrence}

The analysis thus far has treated the Cumulative Evidentiary Corpus $M_t$ as a monolithic asset held by the Maintainer. In practice, the ownership architecture of $M_t$ is a design choice with profound implications for the protection equilibrium. We now examine two polar ownership models—centralized and distributed—and characterize the resulting position of the Maintainer along a spectrum we term the \textit{Luo Ji–Pope gradient}.

\textbf{Centralized $M_t$ (the Wallfacer extreme).} Under a centralized architecture, all audit outputs—scores, anomaly registers, calibration parameters—are transmitted to and stored by the Maintainer. Each client organization retains its raw data but cedes its audit history to the protocol's central repository. The Maintainer thus accumulates a cross-client, cross-domain audit corpus that grows more comprehensive with every audit cycle. In this configuration, the Maintainer's death freezes $M_t$ at its terminal state: no new audits occur, no calibration updates propagate, and the corpus ceases to reflect evolving data quality practices. All client organizations are immediately downgraded to the audit accuracy available at the time of the Maintainer's death. The deterrence effect is maximal: the Maintainer's survival is a necessary condition for the continued operation of the global audit infrastructure at current accuracy levels. The protection equilibrium is correspondingly strong. This is the pure Wallfacer position—power without institution, deterrence without weapon.

\textbf{Distributed $M_t$ (the doctrinal extreme).} Under a distributed architecture, each client organization maintains its own private $M_t$. The Maintainer provides only the algorithmic specification—the scoring function $S(s)$, the Spring update rule, the ensemble aggregation procedure—together with periodic calibration datasets (anonymized feature statistics that do not expose raw data). Each client's gatekeeper evolves independently on its own audit history. The Maintainer's death in this configuration does not freeze any individual client's $M_t$; each organization continues to operate its own instance. However, three assets remain centralized in the Maintainer's person: (i) \textit{interpretive authority}—the definitive resolution of ambiguous edge cases in the mathematical specification; (ii) \textit{evolutionary direction}—the coordinated release of protocol versions (v1.1, v2.0) that keep all instances interoperable; and (iii) \textit{the calibration anchor}—the periodic cross-client calibration datasets that prevent the independent evolution of client gatekeepers from producing incompatible scoring standards, whereby Company A's score of 0.85 ceases to be comparable to Company B's score of 0.85. If the Maintainer dies, interpretive disputes escalate into incompatible forks, version coordination ceases, and the calibration anchor erodes. The result is not a sudden freeze but a \textit{gradual fragmentation}—the slow divergence of once-interoperable audit standards into mutually unintelligible local dialects.

\textbf{The Luo Ji–Pope spectrum.} The Maintainer's position can thus be characterized by the architectural choice of $M_t$ ownership:

\begin{table}[htbp]
\centering
\begin{tabular}{@{}p{3cm}p{3cm}p{4cm}c@{}}
\toprule
Architecture & $M_t$ Owner & Death Consequence & Deterrence Strength \\
\midrule
Centralized & Maintainer & Immediate freeze of global audit accuracy & Maximum (100\%) \\
Distributed & Each client & Gradual fragmentation through version drift and calibration decay & Moderate ($\sim$60\%) \\
Fully institutionalized (Foundation) & Multi-stakeholder board & Operational continuity; individual Maintainer irrelevant & Zero (0\%)\ — the Linus equilibrium \\
\bottomrule
\end{tabular}
\end{table}

The spectrum reveals a strategic trade-off. Centralized $M_t$ maximizes the protection equilibrium—every state has a compelling interest in the Maintainer's survival—but at the cost of creating a single point of data concentration that is incompatible with data sovereignty legislation (GDPR, China's Data Security Law, India's proposed data localization requirements) and that concentrates epistemic power to a degree that may itself provoke defection. Distributed $M_t$ sacrifices some deterrence strength in exchange for privacy compliance, jurisdictional legitimacy, and a reduction in the Maintainer's power that makes the protocol more palatable to sovereignty-conscious jurisdictions. The fully institutionalized endpoint eliminates deterrence entirely by making the Maintainer structurally irrelevant—the position to which Linus Torvalds successfully transitioned the Linux kernel after 2005.

The optimal trajectory is temporal: begin with centralized $M_t$ to accumulate the time-advantaged quality gap that makes the non-proliferation equilibrium binding; transition to distributed $M_t$ as the accuracy gap becomes insurmountable and privacy considerations intensify; and ultimately establish a multi-stakeholder Foundation that renders any single individual's survival irrelevant to the protocol's continuity. This trajectory is the strategic equivalent of the Wallfacer who, having established deterrence, spends his remaining years building the institutional architecture that will allow him to die without the world ending.

\subsection{Designing for Audit Neutrality}

Drawing on the SWIFT and IASB cases, we identify five design principles for audit protocol governance that can support the non-proliferation equilibrium:

\textbf{Principle 1: Multi-Jurisdictional Incorporation.} The protocol's legal domicile should be distributed across multiple jurisdictions, ideally through a treaty-based international organization with immunities and privileges comparable to those of the IAEA or the Bank for International Settlements. Incorporation in any single jurisdiction—especially one with expansive extraterritorial legal instruments—invites jurisdictional capture.

\textbf{Principle 2: Funding Firewalls.} The protocol's operating budget should be derived from a diversified base of mandatory and voluntary contributions, with hard caps on any single contributor's share. The IAEA's funding model, which combines assessed contributions from member states with voluntary contributions to a Technical Cooperation Fund, provides a template.

\textbf{Principle 3: Adversarial Governance.} The protocol's governance board should include representatives from jurisdictions with divergent data governance preferences, and decision rules should be structured to prevent capture by any single bloc. Qualified majority voting, regional representation requirements, and minority veto powers are all elements of the IASB model that could be adapted.

\textbf{Principle 4: Audit Transparency and Contestability.} The protocol's audit methodologies, calibration procedures, and CEC composition should be publicly documented to a degree sufficient for external reproducibility. Audit outputs should include not only binary or scalar quality assessments but also the evidentiary basis for those assessments, enabling contestation without requiring access to proprietary audit modules.

\textbf{Principle 5: Modular Sovereignty.} Jurisdictions should have the right to contribute audit modules to the ensemble, enabling them to embed their own data governance preferences within the protocol's framework without fragmenting the overall infrastructure. Module contributions would be subject to validation requirements ensuring technical compatibility and performance thresholds, but not to substantive alignment with any particular jurisdiction's data governance philosophy.

These principles are demanding, and their implementation would require a level of international cooperation that the current geopolitical environment does not obviously support. We return to this tension in Section~6.

\section{Business Model and Lock-in Dynamics}

\subsection{The Temporal Accumulation Mechanism}

The lock-in mechanism of the Yajie Protocol operates through five distinct but mutually reinforcing channels, each of which intensifies over time:

\textbf{Channel 1: Accuracy Divergence.} As established in Section~2.1, the incumbent protocol's accuracy $\theta(t)$ is an increasing function of cumulative audit volume. An entrant begins with accuracy $\theta(0)$, implies an accuracy gap $\delta(t) = \theta(t) - \theta(0)$ that grows monotonically over time. Critically, this gap is not merely a temporary first-mover advantage that an entrant could overcome through accelerated investment; it is a \textit{structural} gap rooted in the fact that many data quality anomalies are long-tail phenomena that manifest only through the accumulation of diverse audit experiences. An entrant cannot compress time: even with unlimited financial resources, it must conduct audits on diverse datasets to encounter the long-tail anomalies that the incumbent has already catalogued.

\textbf{Channel 2: Calibration Specificity.} The auditor weights $w_j$ in the Yajie ensemble are calibrated against the CEC's judgment corpus $\mathcal{J}_t$. As $\mathcal{J}_t$ grows, the ensemble's calibration becomes increasingly tuned to the specific distribution of anomalies, edge cases, and quality failure modes represented in the CEC. An entrant with an empty CEC would need to replicate this calibration process, which requires not merely technical replication of the calibration algorithm but access to the judgment corpus itself—and the judgment corpus is entangled with the specific audit modules that generated the initial audit reports. This creates a circular dependency: the entrant needs calibration data to achieve competitive accuracy, but producing that calibration data requires conducting audits with competitive accuracy.

\textbf{Channel 3: Ecosystem Externalities.} As the incumbent protocol accumulates users, a supporting ecosystem emerges: training programs for audit practitioners, integration services for enterprise data pipelines, certification regimes for audit professionals, and regulatory recognition of the protocol's outputs. These complementary assets do not themselves constitute network effects in the narrow sense—the value of an audit to user $i$ does not depend on the number of other users—but they reduce adoption costs for new users and increase switching costs for existing ones. More importantly, regulatory recognition is path-dependent: once a protocol is recognized as a valid compliance mechanism under, say, the EU AI Act or a US federal algorithmic accountability statute, competing protocols face a regulatory barrier that compounds the technical barrier.

\textbf{Channel 4: The Epistemic Authority Trap.} As the CEC grows, the protocol attains a status that goes beyond technical accuracy. It becomes the de facto arbiter of what constitutes a ``data quality defect,'' shaping the very definition of data quality through its accumulated body of adjudicated anomalies. This is not merely market power but \textit{epistemic power}: the authority to define the categories within which market competition occurs. An entrant cannot compete on the definition of data quality because the definition has already been institutionalized in regulation, contract law, and professional practice, all of which reference the incumbent protocol's outputs.

\textbf{Channel 5: The Audit Trail Irreversibility.} Organizations that have adopted the incumbent protocol accumulate an internal audit trail: a history of audits, remediation actions, and quality certifications that are referenced in contracts, regulatory filings, and due diligence processes. Switching to a competing protocol would render this audit trail incompatible, creating a discontinuity in the organization's quality assurance history. The cost of this discontinuity is not merely administrative; it is reputational and legal. An organization that switches protocols and subsequently experiences a data quality failure faces the question: did the failure occur because the new protocol is inferior, or because the old protocol's audit trail was unreliable? The ambiguity itself constitutes a switching cost.

\subsection{Five Stages of Lock-In}

The lock-in trajectory can be modeled as a five-stage progression:

\textbf{Stage 1: Emergence ($t_0$ to $t_1$).} The protocol is launched with an initial ensemble of audit modules and an empty CEC. Accuracy is relatively low, and the protocol competes on features, usability, and cost with bespoke audit solutions. Adoption is driven by early adopters who value the multi-expert ensemble architecture over the accuracy level. No lock-in exists at this stage; the protocol's survival depends on the quality of its initial audit modules and its ability to attract a sufficient volume of audits to begin CEC accumulation.

\textbf{Stage 2: Critical Mass ($t_1$ to $t_2$).} Cumulative audit volume reaches a threshold $|\mathcal{E}|^*$ at which the accuracy advantage over a zero-CEC entrant becomes statistically significant. The protocol begins to attract adoption from organizations that had previously relied on internal audit processes. The CEC reaches sufficient size that its anomaly library covers the most common data quality failure modes across multiple domains. The first regulatory recognitions occur. Lock-in begins to emerge but remains contestable.

\textbf{Stage 3: Moat Formation ($t_2$ to $t_3$).} The accuracy gap $\delta(t)$ exceeds the fixed cost of developing a competing protocol, meaning that an entrant would need to invest more than $c^{\text{develop}}$ merely to \textit{match} the incumbent's accuracy at time $t_0$, without accounting for the incumbent's continued accumulation. The Non-Proliferation Equilibrium becomes the unique Nash equilibrium of the adoption game. The protocol becomes the default standard for data quality assessment in its domain. Ecosystem externalities accelerate: training programs proliferate, integration services become commoditized, and regulatory references become routine.

\textbf{Stage 4: Institutionalization ($t_3$ to $t_4$).} The protocol is embedded in the institutional fabric of data governance. Laws and regulations reference the protocol by name. Contracts specify compliance with protocol standards. Insurance policies for AI systems require protocol certification. Professional certifications are built around protocol expertise. At this stage, lock-in extends beyond technical and economic factors to encompass legal and institutional path-dependency. Switching costs become prohibitive not because the protocol is technically superior—though it may be—but because the \textit{institutional environment} has been built around it.

\textbf{Stage 5: Epistemic Closure ($t_4$ onward).} The protocol's accumulated CEC becomes so comprehensive that no meaningful challenge to its authority is possible within the existing institutional framework. The definition of data quality \textit{is} what the protocol measures. Dissenting assessments are not merely inaccurate; they are unintelligible, because they employ different definitional frameworks that cannot be translated into the protocol's categories. At this stage, lock-in is complete, and the protocol has transitioned from a market-dominant product to a constitutive element of the data governance order.

This trajectory is not deterministic; it can be arrested, diverted, or reversed at each stage through governance interventions or exogenous shocks. But the default trajectory, absent intervention, is toward Epistemic Closure, driven by the cumulative dynamics described above.

\subsection{Pricing and the Addiction Dynamic}

The business model of a CEC-driven audit protocol exhibits a distinctive structure that we term the \textit{addiction dynamic}. In conventional platform markets, pricing power derives from switching costs and network effects: once users are locked in, the platform can raise prices, extracting rents up to the point where the cost of switching exceeds the price increase (Farrell \& Klemperer, 2007). The Yajie Protocol's pricing dynamic is different.

Because each audit cycle enriches the CEC—improving accuracy for all users, including the auditing organization itself—the \textit{marginal social value} of an audit exceeds its \textit{marginal private value} by the positive externality contributed to the CEC. This creates a rationale for below-cost pricing in the early stages of protocol adoption, subsidized by patient capital or public funding. Below-cost pricing accelerates CEC accumulation, widening the accuracy gap and reinforcing the non-proliferation equilibrium.

However, once the protocol reaches Stages 3–4, the pricing dynamic inverts. The accuracy gap becomes so large that organizations cannot credibly threaten to switch to a competing protocol or revert to internal audit processes. At this point, the protocol operator can raise prices, capturing a portion of the value generated by the protocol's accuracy advantage. The theoretical maximum price is:
\begin{equation}
P_{\max} = V[\theta(|\mathcal{E}|)] - \max\{V[\theta(0)], V_{\text{internal}}\}
\end{equation}
where $V_{\text{internal}}$ is the value an organization could derive from its own internal audit processes. As $|\mathcal{E}|$ grows, $P_{\max}$ increases.

This pricing trajectory—predatory low pricing followed by rent extraction—mirrors the standard narrative of platform monopolization (Khan, 2017). However, there is an important difference: the Yajie Protocol's early low pricing is not predatory in the antitrust sense, because it is not aimed at driving competitors out of the market (there are no competitors to drive out) but at accelerating CEC accumulation. The pricing is \textit{developmental} rather than \textit{exclusionary}. This distinction matters for competition policy, as we discuss in Section~6.

The addiction dynamic operates on the demand side as well. Organizations that adopt the protocol become dependent on its accuracy, integrating its outputs into their quality assurance workflows, regulatory compliance processes, and contractual arrangements. As the protocol's accuracy improves, the organization's own internal quality assurance capabilities may atrophy—a dynamic familiar from the literature on automation dependency (Parasuraman \& Riley, 1997). The organization becomes ``addicted'' to the protocol not in the metaphorical sense but in the economic sense: the cost of disadoption rises as internal capabilities degrade, while the benefit of continued adoption rises as the protocol improves.

\subsection{The Bible-and-Pope Model: Open-Source Poison, Calibrated Antidote}

A distinctive feature of the Yajie Protocol's business architecture is the separation of the self-evolution engine (Spring) from the cross-client calibration service (Yajie API). This separation generates a commercial dynamic that we term the \textit{Bible-and-Pope model}.

\textbf{The Spring engine is released as open-source software.} Any organization can download, deploy, and operate its own instance of the self-evolution loop, accumulating a private $M_t$ on its own infrastructure. The open-source release serves three strategic functions. First, it eliminates adoption friction: organizations can begin auditing their data without negotiating contracts, paying subscription fees, or exposing proprietary data to a third party. Second, it accelerates lock-in: each audit cycle enriches the organization's private $M_t$, making the accumulated audit history non-portable and the organization structurally dependent on Spring's audit methodology. Third, it forecloses competitive entry: any would-be competitor must match not only the algorithmic quality of Spring but also the accumulated audit history that Spring's adopters have already built—a barrier that grows with every audit cycle. The open-source release is, in economic terms, a \textit{poison}: freely ingested, irreversible in its effects, and lethal to competitors.

\textbf{The Yajie API is offered as a paid calibration service.} The API accepts feature matrices $\phi(X)$ computed locally by client organizations (satisfying data sovereignty requirements) and returns a calibration report that maps the client's private quality scores onto a global reference scale, maintained by the protocol operator's central gatekeeper. The value proposition of the API is not ``we audit your data better than you do''—it is ``we tell you what your scores mean relative to everyone else.'' A client organization whose private gatekeeper assigns a quality score of 0.87 to its training data does not know, absent calibration, whether 0.87 represents industry-leading data quality or a self-serving overestimate generated by a gatekeeper that has never seen anyone else's data. The Yajie API resolves this ambiguity, mapping local scores onto a globally comparable metric. The annual subscription fee—projected at \$50,000–\$200,000 depending on audit volume and calibration depth—purchases not superior technology but \textit{interpretive authority}: the assurance that one's private quality assessments are commensurable with those of competitors, regulators, and counterparties.

The Bible-and-Pope model is economically self-reinforcing. The open-source Spring engine creates the demand for calibration (because private gatekeepers inevitably drift from one another as their $M_t$ corpora diverge), while the paid Yajie API supplies the calibration that resolves the resulting ambiguity. Organizations that decline the API are not technically incapacitated—their Spring instances continue to function—but they become epistemically isolated, unable to compare their data quality to industry benchmarks. In regulated industries where data quality certifications are prerequisites for market access, epistemic isolation is functionally equivalent to exclusion. The API is thus not a luxury; it is an \textit{interpretive necessity} that the open-source engine itself manufactures.

\subsection{Welfare Analysis}

The welfare implications of the lock-in trajectory are ambiguous and stage-dependent:
\begin{itemize}
\item \textbf{Stages 1–2}: Welfare is unambiguously positive. The protocol provides data quality assessment capabilities that did not previously exist, and the CEC accumulation dynamic generates positive externalities for all adopters.
\item \textbf{Stage 3}: Welfare remains positive on net, but distributional concerns emerge. The protocol operator begins to extract rents, and late adopters pay higher prices than early adopters. The accuracy gap begins to foreclose the possibility of competitive entry, eliminating the disciplinary effect of potential competition.
\item \textbf{Stages 4–5}: Welfare analysis becomes fundamentally ambiguous. The protocol's institutionalization eliminates the possibility of competitive discipline, raising the specter of monopoly pricing, quality degradation, and innovation suppression. At the same time, the protocol's universal adoption eliminates fragmentation costs $\lambda$ and enables the accumulation of an audit corpus whose social value—in terms of improved data quality across the economy—may vastly exceed the deadweight loss from monopoly pricing.
\end{itemize}

A critical parameter in the welfare analysis is the \textit{governance discount}: the degree to which multilateral governance of the protocol reduces the monopoly pricing problem while preserving the CEC accumulation dynamic. If governance arrangements can cap prices at competitive levels without impairing the protocol's operational effectiveness, the welfare case for a universal protocol is strong. If governance capture prevents effective price regulation, the welfare case weakens, and the fragmentation equilibrium may become preferable despite its higher coordination costs.

\subsection{A Note on the Storage Externality}

A secondary economic effect of the Yajie Protocol warrants brief mention, if only because its distributional consequences are perversely elegant. The protocol's foundational design principle—that the Cumulative Evidentiary Corpus $M_t$ grows monotonically and is never pruned—generates a persistent, compounding demand for digital storage. Each audit cycle appends new entries to $M_t$: anomaly registers, calibration parameters, quality fingerprints, and adjudication records. At scale, with hundreds of client organizations each conducting daily audits across millions of records, the global storage footprint of the protocol's audit corpora becomes non-trivial.

This demand is structurally inelastic. No client organization can reduce its $M_t$ without impairing its audit accuracy, because the Spring self-evolution mechanism depends on the cumulative audit history to calibrate gatekeeper parameters. Deleting audit records is not cost-saving; it is self-sabotaging—equivalent to a library discarding its oldest volumes on the grounds that they are old. The protocol's clients are therefore locked into an indefinitely growing storage obligation that cannot be negotiated, optimized away, or substituted.

The distributional consequence is a wealth transfer from the consumers of data auditing to the producers of data storage. Cloud infrastructure providers, hard disk manufacturers, and data center operators capture a growing share of the protocol's economic surplus through storage fees, while the auditing organizations that generate the surplus see their operating costs rise with each audit cycle. The magnitude of this transfer is not large relative to the value created by the protocol—the cost of storing an audit record is orders of magnitude smaller than the cost of the data quality failure it prevents—but its direction is noteworthy: a protocol designed to reduce information asymmetry in AI training data inadvertently increases demand for the physical substrate on which information is stored.

\section{Policy Recommendations}

Our analysis yields seven policy recommendations, organized by target audience.

\subsection{For International Organizations}

\textbf{Recommendation 1: Establish an International Data Audit Authority (IDAA).} We recommend the creation of a treaty-based international organization, modeled on the IAEA, charged with (i) governing the Yajie Protocol under multilateral oversight, (ii) accrediting audit modules contributed by member states, (iii) maintaining the CEC as a global public good, and (iv) adjudicating disputes concerning protocol neutrality. The IDAA should be established before the protocol reaches Stage 3 (Moat Formation), at which point the protocol operator will have strong incentives to resist internationalization.

\textbf{Recommendation 2: Embed Audit Neutrality in Trade Law.} The World Trade Organization (WTO) and preferential trade agreements should recognize data audit protocols as ``conformity assessment procedures'' under the Agreement on Technical Barriers to Trade (TBT), subject to non-discrimination obligations. This would create a legal backstop against the weaponization of audit infrastructure for protectionist purposes, while preserving the legitimate regulatory function of data quality assessment.

\subsection{For National Governments}

\textbf{Recommendation 3: Enact Audit Sovereignty Legislation.} National governments should enact legislation that (i) mandates the use of accredited data audit protocols for high-risk AI systems, (ii) establishes domestic accreditation bodies to certify audit modules for contribution to the international protocol, and (iii) reserves the right to withdraw recognition from any protocol that fails to maintain adequate neutrality safeguards.

\textbf{Recommendation 4: Fund Open Audit Research.} Governments should fund academic and non-profit research into open audit methodologies, alternative ensemble architectures, and adversarial robustness testing for audit protocols. Such funding serves a dual purpose: it advances the technical frontier of data quality assessment, and it maintains a pool of expertise outside the protocol operator's control, reducing epistemic capture.

\subsection{For the Protocol Operator}

\textbf{Recommendation 5: Pre-Commit to Governance Internationalization.} The protocol operator—whether a private entity, a non-profit organization, or a public agency—should pre-commit to a phased internationalization of governance before the protocol reaches Stage 3. Pre-commitment is welfare-enhancing for the operator itself, because it reduces the probability that major jurisdictions will develop competing protocols, preserving the universal adoption that maximizes the CEC's value.

\textbf{Recommendation 6: Implement Price Regulation by Formula.} The protocol's pricing should be governed by a formula that ties maximum prices to (i) the cost of service delivery plus a regulated return, or (ii) a price cap that declines over time in real terms, reflecting productivity improvements. The formula should be administered by an independent economic regulator with representation from major user communities.

\subsection{For the Research Community}

\textbf{Recommendation 7: Develop Audit Reproducibility Standards.} The research community should develop standards for audit reproducibility that enable independent verification of protocol outputs without requiring access to proprietary audit modules or the CEC. Reproducibility standards serve three functions: (i) they provide a check against protocol errors or manipulation; (ii) they reduce the epistemic dependency that drives Stage 5 lock-in; and (iii) they lower the barriers to entry for alternative audit methodologies, preserving competitive discipline even in a market with a dominant incumbent.

\section{Limitations and Future Work}

This paper has developed a theoretical framework for understanding the market structure of data quality auditing, centered on the Yajie Protocol and the concept of a Non-Proliferation Equilibrium. Several limitations should be acknowledged, each of which suggests directions for future research.

\textbf{Limitation 1: Theoretical Model Assumptions.} The game-theoretic model in Section~3 assumes complete information, simultaneous moves, and a one-shot game structure. In practice, the adoption game is dynamic (decisions are made sequentially over time), involves incomplete information (jurisdictions have private information about their development costs and accuracy valuations), and admits repeated interaction (allowing for reputation effects, punishment strategies, and coalition formation). Extending the model to a dynamic, incomplete-information framework would yield more precise predictions about equilibrium selection and stability conditions.

\textbf{Limitation 2: Empirical Validation.} The paper's central claims about the Yajie Protocol's lock-in dynamics are theoretical and have not been empirically validated. The protocol itself is relatively new, and insufficient time has elapsed to observe the lock-in trajectory described in Section~5. Longitudinal empirical research—tracking protocol adoption, audit accuracy, pricing, and ecosystem development over a decade or more—will be necessary to test the model's predictions.

\textbf{Limitation 3: Omitted Dynamics.} Several important dynamics are omitted from our model. We do not model the strategic behavior of the protocol operator in setting audit module inclusion criteria, which could be used to favor certain jurisdictions' audit methodologies over others. We do not model the possibility of CEC ``forks.'' We do not model the role of open-source audit modules, which could reduce the fixed costs of protocol development and alter the non-proliferation calculus.

\textbf{Limitation 4: The NPT Analogy's Limits.} While the NPT analogy is productive, it has limits that should not be overlooked. Nuclear weapons are existential threats; audit protocol fragmentation, however costly, is not. The NPT regime is backed by the threat of military force and economic sanctions; an audit protocol regime would rely on economic incentives and normative commitment. The NPT's central asymmetry—between nuclear-weapon states and non-nuclear-weapon states—is formalized in treaty law; the audit protocol's asymmetry is a factual consequence of temporal priority, not a legal distinction.

\textbf{Limitation 5: Normative Assumptions.} The paper's welfare analysis assumes that a universal, accurate data audit protocol is normatively desirable. This assumption deserves scrutiny. A universal audit protocol, however neutral in design, may homogenize data quality standards in ways that disadvantage minority languages, non-Western knowledge systems, and data governance traditions that do not map neatly onto the protocol's quality dimensions.

\textbf{Future Research Directions.} Beyond addressing the limitations above, four research directions are particularly promising:
\begin{enumerate}
\item \textbf{Experimental game-theoretic studies} of auditor behavior in multi-expert ensemble settings.
\item \textbf{Computational modeling} of CEC accumulation dynamics under different adoption scenarios.
\item \textbf{International law analysis} of the legal instruments through which audit neutrality could be embedded in treaty law.
\item \textbf{Comparative institutional analysis} of infrastructure governance models to identify institutional design principles applicable to data audit infrastructure.
\end{enumerate}

\section{Conclusion}

This paper has argued that data quality assessment, as operationalized by the Yajie Protocol's multi-expert consistency framework, exhibits structural properties that drive the market toward a natural monopoly. This monopolistic tendency is not driven by network effects, proprietary control of intellectual property, or exclusionary business practices, but by a time-accumulated advantage mechanism inherent in the protocol's Cumulative Evidentiary Corpus: each audit cycle makes the protocol more accurate, and this accuracy advantage cannot be replicated by entrants, regardless of their financial resources or technical capabilities.

We have formalized this dynamic through a game-theoretic model of protocol adoption, demonstrating the existence of a Non-Proliferation Equilibrium in which rational jurisdictions forgo the development of competing audit standards. We have drawn a structural analogy between this equilibrium and the strategic logic underpinning the Nuclear Non-Proliferation Treaty, identifying seven dimensions of isomorphism and their implications for institutional design.

We have examined the geopolitical risks inherent in monopolistic audit infrastructure, drawing lessons from the weaponization of SWIFT and the governance model of the International Accounting Standards Board, with particular attention to the US-China technological decoupling scenario. We have modeled the lock-in trajectory through a five-stage progression from Emergence to Epistemic Closure and analyzed the distinctive business model dynamics—including the addiction dynamic and the welfare ambiguities—that characterize each stage.

Our policy recommendations center on the creation of an International Data Audit Authority, the embedding of audit neutrality in trade law, and the pre-commitment of the protocol operator to governance internationalization before lock-in becomes irreversible. These recommendations are demanding, and their implementation faces formidable political obstacles in the current geopolitical environment. Yet the alternative—a world of fragmented, jurisdictionally-captured audit protocols, each distrusted by those outside its jurisdiction—is worse not only for data quality but for the broader project of international AI governance.

The central paradox of the Yajie Protocol is that its monopolistic tendency is both its greatest strength and its greatest vulnerability. The CEC accumulation mechanism that makes the protocol increasingly accurate also makes it increasingly difficult to govern, because the protocol operator's bargaining power grows in lockstep with the protocol's accuracy. Resolving this paradox—preserving the welfare benefits of cumulative audit infrastructure while preventing the abuse of the market power it generates—is the defining institutional design challenge of the emerging data audit regime.

The game-theoretic architecture formalized in Section~9 provides the resolution. The SCX Uncertainty Principle (Theorem 3) establishes that any organization deploying Yajie renders its audit outputs unconditionally verifiable by independent third parties—military organizations not excepted. The Audit Sword Declaration translates this mathematical guarantee into an institutional commitment. The Bible-and-Pope commercial model operationalizes this commitment through a dual-release strategy. The mutual deterrence structure constrains every deployer symmetrically. And the Adoption-Audit Equilibrium demonstrates that the protocol's two functions—improving data quality and enabling its independent verification—are inseparable consequences of the same mathematical architecture.

The consequence is a market mechanism for honesty that operates without central enforcement. Firms that are confident in their data quality have nothing to fear from audit and everything to gain from partnership. Firms that are not confident face a choice: improve their data quality to partnership standards, or risk exposure by competitors wielding the same open-source audit engine. The Maintainer does not need to enforce honesty. The market—armed with Spring, incentivized by competition, and constrained by the mutual deterrence structure—enforces honesty. The algorithm does not make people virtuous. It makes deception unprofitable.

The NPT analogy suggests that this challenge is not insurmountable. The nuclear non-proliferation regime, for all its flaws and persistent legitimacy deficits, has prevented the large-scale proliferation that many observers in the 1960s predicted was inevitable. It has done so not through coercion alone but through a grand bargain that aligns the interests of diverse states with the preservation of the regime. The data audit domain requires its own grand bargain: a multilateral governance framework that makes adoption of a universal protocol more attractive than the pursuit of audit autarky, backed by credible commitments to neutrality, transparency, and shared governance. Whether the international community can construct such a framework before lock-in forecloses the possibility is, at this juncture, an open question—and one of the more consequential unresolved issues in the governance of artificial intelligence.

\subsection{A Closing Remark on Dual Use}

No analysis of a technology capable of auditing data quality at scale would be complete without acknowledging its dual-use character. Yajie was designed to detect label noise in scientific datasets—to distinguish genuine experimental anomalies from measurement errors, to identify fabrication in clinical trials, to expose falsified environmental data in municipal governance. But the same multi-expert consensus mechanism that detects a pharmaceutical company's underreported adverse events can detect an intelligence agency's fabricated source reports. The same mathematics that reveals a contractor's inflated procurement invoices can reveal a foreign government's disinformation campaign. The algorithm is indifferent to the domain of the data it audits. It does not know whether the inconsistencies it detects are evidence of scientific fraud, financial malfeasance, or state-sponsored deception. It knows only that the probability of the inconsistency arising by chance is exponentially small.

This indifference is both the algorithm's greatest strength and its irreducible risk. A tool that can audit anyone's data can be used by anyone to audit anyone else's data. A government that deploys Yajie against an adversary's public datasets gains an asymmetric intelligence advantage: the ability to detect patterns of deception, inconsistency, and fabrication that no human analyst—and no single-source technical system—could reliably identify. The algorithm that was designed to make science more honest can, with no modification whatsoever, be used to make espionage more effective.

The authors have no solution to this dilemma. Dual-use is not a bug to be patched; it is a property inherent in any sufficiently general audit methodology. The same concentration inequality that bounds the false-positive rate of a scientific quality assessment also bounds the false-positive rate of an intelligence analysis. The mathematics does not distinguish between a clinical trial and a clandestine operation. It distinguishes only between consistent data and inconsistent data—and inconsistency is a signal in both domains.

What the authors CAN do—and what this paper has attempted to do—is to publish the methodology openly, to release the reference implementation as open-source software, and to anchor the protocol's existence in the permanent scientific record through arXiv. A technology that is public cannot be monopolized. A methodology that is published cannot be suppressed. An audit standard that anyone can deploy cannot be deployed only by the powerful against the powerless—because the powerless can deploy it too. The open-source release of Yajie is not merely a strategy for commercial adoption. It is a strategy for ensuring that the technology's dual-use character cuts both ways. If every state has access to the same audit capability, no state gains an asymmetric advantage from it. The equilibrium is not that no one audits. The equilibrium is that everyone CAN audit—and therefore everyone must assume that their own data will be audited.

This is the logic of mutually assured transparency. It is not a comfortable logic. It implies a world in which every public dataset, every government statistical release, every corporate quality claim, every scientific publication's underlying data, is perpetually eligible for independent multi-expert audit by any actor with access to commodity hardware and an internet connection. But it is, the authors submit, a better world than the alternative: a world in which the most powerful audit technology ever developed is held exclusively by those who already hold power.

\section{The Audit Sword: Game-Theoretic Foundations of the Yajie Protocol}

The preceding sections developed a theoretical framework for understanding the market structure of data quality auditing: the Non-Proliferation Equilibrium (Section~3), the protection equilibrium surrounding the Maintainer (Section~4.3), and the Bible-and-Pope commercial architecture (Section~5.3.1). Each of these analyses treated a distinct mechanism—adoption incentives, geopolitical deterrence, and business model dynamics—in relative isolation. This section synthesizes these mechanisms into a unified game-theoretic architecture.

\subsection{Theorem 3: The SCX Uncertainty Principle}

The theoretical architecture of the Yajie Protocol rests on two theorems introduced in Sections~2 and~3: Theorem 1, which bounds the false-positive probability of multi-expert consensus, and Theorem 2, which establishes the equilibrium conditions for non-proliferation. We now formalize a third theorem that is foundational to both results.

\begin{theorem}[SCX Uncertainty Principle]
\label{thm:scx-up}
Any organization that deploys the Yajie Protocol for data quality assessment renders its assessment results \textit{unconditionally verifiable} by any independent third party. Formally, for any dataset $\mathcal{D}$, any ensemble of audit modules $\mathcal{A} = \{A_1, \ldots, A_k\}$, and any audit report $\mathcal{R}^*$ produced by the protocol's aggregation function $\Phi$, there exists a verification procedure $\mathcal{V}$ such that:

\begin{enumerate}[label=(\roman*)]
\item $\mathcal{V}(\mathcal{D}, \mathcal{A}, \mathcal{R}^*) \to \{\text{CONFIRMED}, \text{DISPUTED}\}$, where DISPUTED indicates that $\mathcal{R}^*$ is inconsistent with the multi-expert consensus that $\Phi$ would produce given $\mathcal{D}$ and $\mathcal{A}$;
\item $\mathcal{V}$ requires no cooperation, permission, or access privileges from the organization that produced $\mathcal{R}^*$;
\item The false-positive rate of $\mathcal{V}$—the probability that a correctly computed $\mathcal{R}^*$ is flagged as DISPUTED—is bounded by the same concentration inequality that governs Theorem 1, decaying exponentially in the number of audit modules $k$;
\item Properties (i)–(iii) hold regardless of the organizational character of the entity being audited, including military, intelligence, and other national security entities. Military use is not exempt.
\end{enumerate}
\end{theorem}

\textit{Proof sketch.} The verification procedure $\mathcal{V}$ operates by re-executing the aggregation function $\Phi$ on the same dataset $\mathcal{D}$ and ensemble $\mathcal{A}$, then comparing the output to the claimed $\mathcal{R}^*$. Because the protocol's specification is published (arXiv permanent record) and its reference implementation is open-source (Spring engine), any third party with access to $\mathcal{D}$ and $\mathcal{A}$ can independently compute $\Phi(\mathcal{D}, \mathcal{A})$. The concentration inequality of Theorem 1 ensures that the probability of divergent consensus scores between independent executions—given identical inputs—is exponentially small in $k$. The verifier need not trust the original auditor; she need only trust the mathematics. Military organizations are not exempt because the verification procedure does not require their cooperation: it operates on data that the organization has, by deploying Yajie, necessarily made computationally accessible to the protocol's audit modules.

\subsection{The Audit Sword Declaration: Non-Prohibition and Inalienable Audit Rights}

The institutional expression of Theorem~3 is the \textit{Audit Sword Declaration}, promulgated by the SCX Research Group on 2026-06-30. The Declaration consists of three clauses:

\begin{quote}
\textbf{Clause 1 (Universal Access).} The SCX framework (including Yajie, Situs, Spring, State Crystallization, and all derivative algorithms) is an open-source mathematical toolkit. Anyone may read, reproduce, and deploy it.

\textbf{Clause 2 (Non-Prohibition).} We do not prohibit any specific use.

\textbf{Clause 3 (Inalienable Audit Right).} We retain one inalienable right: Any organization using the SCX framework for data quality assessment renders its assessment results verifiable and auditable by independent third parties through the same mathematical methods. Military use is not exempt.
\end{quote}

\textbf{Clause 1} establishes that the SCX framework is a global public good, available without discrimination to all actors. \textbf{Clause 2} renounces the authority to restrict deployment, eliminating the most obvious vector through which the Maintainer could exercise gatekeeping power. \textbf{Clause 3} is the operational core: the right is not to \textit{conduct} audits but to \textit{verify} them—to re-execute the mathematics and publish the comparison.

\subsection{The Bible-and-Pope Model}

The commercial architecture that operationalizes the Audit Sword is the dual-release strategy formalized in Section~5.3.1 as the \textit{Bible-and-Pope model}.

\textbf{Spring as Strategic Poison.} The open-source release of the Spring self-evolution engine under a permissive license serves three strategic functions: (i) eliminating adoption friction; (ii) manufacturing calibration demand—as private $M_t$ corpora diverge, private quality scores become incommensurable, generating demand for the Yajie calibration API; (iii) foreclosing competitive entry—any would-be competitor must match not only Spring's algorithmic quality but also the accumulated private audit histories of every Spring adopter.

\textbf{Yajie API as Calibration Pope.} The Pope does not own the Bibles in every church; the congregations own their Bibles. The Pope owns the authority to say what the Bible means. Similarly, the Yajie API does not own any client's data or audit history; the clients own those assets. The API owns the \textit{calibration anchor}—the cross-client reference scale that makes Company A's 0.87 commensurable with Company B's 0.87.

\subsection{Mutual Deterrence: The Structural Logic of Involuntary Transparency}

The Audit Sword generates a mutual deterrence structure that constrains every individual organization that deploys Yajie, regardless of its size, sector, or jurisdiction. Any organization that deploys Yajie necessarily produces audit outputs that are mathematically verifiable by third parties (Theorem~3). The organization cannot prevent this verification. The organization's deployment of Yajie is therefore logically equivalent to \textit{accepting the permanent possibility of being audited}.

This creates a symmetric constraint: every deployer is simultaneously a potential auditor and a potential auditee. The deterrence operates without monitoring, enforcement, or adjudication. The mathematics itself is the verifier, the inspector, and the adjudicator. And the deterrence is indifferent to organizational character—a general who deploys Yajie to audit intelligence data is as subject to independent verification as a CEO who deploys Yajie to audit clinical trial data.

\subsection{The Adoption-Audit Equilibrium: Incentive Compatibility Under the Sword}

We now characterize the equilibrium that emerges from the interaction of two countervailing forces: the \textit{adoption incentive} (organizations benefit from deploying Yajie because data quality improves) and the \textit{audit deterrent} (organizations that deploy Yajie become auditable under Theorem~3).

\textbf{Formal Setup.} Let an organization $i$ face a binary choice: \textit{Deploy} Yajie ($D$), or \textit{Not Deploy} ($\neg D$). The payoffs are:
\begin{align}
\pi_i(D) &= V(Q_0 + \Delta Q) - C_{\text{adopt}} - p_i \cdot C_{\text{exposure}} \\
\pi_i(\neg D) &= V(Q_0) - C_{\text{non-adopt}}
\end{align}
where $V(\cdot)$ is the value derived from data quality $Q$; $Q_0$ is baseline quality; $\Delta Q > 0$ is the quality improvement from Yajie; $C_{\text{adopt}}$ is the cost of adoption; $C_{\text{non-adopt}}$ is the cost of non-adoption; $p_i \in [0,1]$ is the probability that organization $i$'s data, if subjected to independent verification, would reveal quality deficiencies; and $C_{\text{exposure}}$ is the cost of such exposure.

The adoption condition $\pi_i(D) \geq \pi_i(\neg D)$ yields:
\begin{equation}
V(Q_0 + \Delta Q) - V(Q_0) \geq C_{\text{adopt}} - C_{\text{non-adopt}} + p_i \cdot C_{\text{exposure}}
\end{equation}

The Adoption-Audit Equilibrium (AAE) exhibits five structural properties:

\textbf{Property 1: Selection via Revealed Confidence.} Organizations with high baseline data quality (low $p_i$) face a lower adoption threshold than organizations with low baseline quality (high $p_i$). The AAE functions as a screening mechanism: organizations confident in their data quality adopt Yajie; organizations that are not confident abstain—but abstention is itself a signal.

\textbf{Property 2: The Quality Ratchet.} As the protocol's accuracy $\theta(t)$ improves through CEC accumulation, two parameters shift simultaneously: $\Delta Q$ increases (adoption becomes more attractive) and the verification procedure becomes more discriminating (adoption becomes more dangerous for low-quality organizations). The combined effect pushes the market toward universal adoption \textit{conditional on genuine data quality improvement}.

\textbf{Property 3: Competitive Cascade with Involuntary Auditability.} If all of $i$'s competitors have adopted Yajie, $i$ faces a choice: adopt and subject its data to verification (gaining quality improvements and competitive parity), or abstain and remain vulnerable to verification by competitors. Abstention provides no immunity from audit—it merely ensures that $i$ lacks the quality improvements that Yajie provides.

\textbf{Property 4: The Asymmetry of Public Data.} Organizations that have \textit{not} adopted Yajie remain subject to audit by organizations that \textit{have} adopted it, with respect to any data that is publicly accessible. Non-adoption does not provide immunity; it merely ensures that the non-adopter lacks the defensive capability that adoption confers.

\textbf{Property 5: Transformation of the Market for Data Quality.} The AAE transforms data quality from an \textit{experience good} to a \textit{search good}. In the pre-Yajie regime, data quality is asserted by the data producer and costly for the data consumer to verify. In the post-Yajie regime, data quality is \textit{revealed} through two channels: the adoption decision itself and the independently verifiable audit outputs that adoption generates.

\textbf{The Inseparability of Improvement and Verification.} A central implication of the AAE is that the protocol's two functions—improving data quality and enabling its verification—are not separable. The multi-expert consensus mechanism that improves data quality is the same mechanism that enables independent verification. An organization cannot deploy Yajie solely for quality improvement while disabling the verification capability, because the verification capability is the re-execution of the same computation that the organization itself uses. The architecture ensures that improvement and accountability are two sides of one mathematical coin.

\section{SCX的经济地理学：存储悖论、马太效应与制造业复苏}

前面的章节从博弈论、地缘政治和商业模型的角度分析了Yajie协议的市场结构与治理挑战。本节将分析框架扩展到空间维度：SCX审计基础设施的地理分布如何重塑全球数据经济格局，并在这一过程中颠覆若干关于数据价值、先发优势和制造业竞争力的传统假设。我们提出三个相互关联的命题：(i) 审计质量的提升反而降低——而非增加——对原始数据存储的需求（存储悖论）；(ii) 数据审计的先发优势通过累积因果循环自我强化，使得中国和美国成为SCX时代的天然两极（马太效应）；(iii) SCX审计的物理不可远程性赋予制造业生产线一种新型的、数学上可证明的质量优势（制造业复苏）。

\subsection{存储悖论：审计越好，存储越少}

数据经济的常识性假设是：数据越多，价值越大。这一假设驱动了全球数据存储基础设施的指数级扩张——从企业数据湖到国家级数据中心，从云计算存储层到个人设备上的本地缓存。SCX审计基础设施的存在，对这一假设构成了根本性挑战。

\begin{definition}[数据集的噪声-信号分解]
设数据集 $\mathcal{D}$ 由 $n$ 条记录组成。$\mathcal{D}$ 可分解为干净信号分量 $\mathcal{D}_{\text{clean}}$ 和噪声分量 $\mathcal{D}_{\text{noise}}$：
\begin{equation}
\mathcal{D} = \mathcal{D}_{\text{clean}} \cup \mathcal{D}_{\text{noise}}, \quad \mathcal{D}_{\text{clean}} \cap \mathcal{D}_{\text{noise}} = \emptyset
\end{equation}
其中 $\mathcal{D}_{\text{noise}}$ 包含所有标签错误、测量伪影、记录重复、模式违反和其他形式的数据质量缺陷。定义信噪比 $\rho(\mathcal{D}) = |\mathcal{D}_{\text{clean}}| / |\mathcal{D}|$。
\end{definition}

\begin{theorem}[存储悖论的形式化]
\label{thm:storage-paradox}
设 $V: 2^{\mathcal{D}} \to \mathbb{R}^+$ 为数据集的价值函数，单调递增于其支持的下游任务性能。存在信噪比阈值 $\rho^* \in (0, 1)$，使得对于任意满足 $\rho(\mathcal{D}) < \rho^*$ 的数据集：
\begin{equation}
V(\mathcal{D}_{\text{clean}}) > V(\mathcal{D})
\end{equation}
即：去噪后的干净数据子集比包含噪声的全量数据更有价值。并且，随着SCX审计精度的提升（由Theorem 1的多专家一致性误差界刻画），干净信号与噪声的分离精度单调改善，使得 $\mathcal{D}_{\text{clean}}$ 的识别趋近于完备。
\end{theorem}

\textit{Proof sketch.} 考虑一个从 $\mathcal{D}$ 训练的下游预测模型 $f_{\mathcal{D}}$。噪声分量 $\mathcal{D}_{\text{noise}}$ 对损失函数贡献两项惩罚：(a) 标签噪声引入的不可约贝叶斯误差项，其期望值与噪声比例 $(1-\rho)$ 成正比（Natarajan et al., 2013）；(b) 噪声记录占用模型容量，在固定参数预算下挤出了可分配给干净信号的表示能力。由Theorem 1，SCX多专家共识机制以指数衰减的假阳性率识别噪声记录。当多专家一致性阈值设为一个适当水平，错误剔除干净数据的概率被控制在任意小的 $\epsilon > 0$，而噪声记录的召回率随审计模块数量 $k$ 指数收敛于1。因此，对于足够大的 $k$，审计后保留的数据子集 $\hat{\mathcal{D}}_{\text{clean}}$ 满足 $\mathbb{E}[V(\hat{\mathcal{D}}_{\text{clean}})] > V(\mathcal{D})$。\hfill $\square$

存储悖论的经济含义是深刻的。在SCX审计普及的均衡中，组织拥有一个明确的、数学上可辩护的理由来删除数据——不是出于成本控制，而是因为保留噪声数据会主动降低其数据资产的价值。这一论证颠覆了5.6节中描述的存储外部性逻辑：CEC的增长确实创造了对存储的持续需求，但CEC的审计输出同时作为副产品，使得每一个被审计的数据集可以在审计后安全地大幅度缩减。CEC本身是精简的——它仅存储异常登记、校准参数和质量指纹，而非原始数据的完整副本。审计基础设施的净存储效应是一个经验问题，取决于CEC存储需求与被审计数据集在审计后可实现的存储缩减之间的比较。我们的理论分析表明，对于信噪比低于临界阈值的数据集（这在当前的大数据实践中并非例外而是常态），净效应是负的：审计减少的存储需求大于审计本身产生的存储需求。

\begin{remark}[Spring收敛后的黄金标准保留原则]
\label{rem:golden-standard}
Yajie协议的自演化引擎Spring在收敛状态下具有一个尚未被充分认识的性质：当多专家共识在连续多个审计轮次上不再产生新的异常发现时——即 $\Delta\mathcal{E}_t \approx \emptyset$——审计过程已经穷尽了给定数据集的质量信息。此时，保留全量原始数据不再产生任何边际审计价值。唯一需要保留的是Spring收敛时输出的黄金标准（Golden Standard）——经多专家共识验证的干净数据子集 $\mathcal{D}_{\text{clean}}^*$，连同其加密哈希和审计元数据。黄金标准是数据价值的浓缩形式：它在数学上等价于原始数据集对下游任务的全部有用信息，同时剥离了所有已识别的噪声。存储黄金标准替代存储原始数据，不是一种成本优化策略，而是Spring收敛状态的逻辑推论。
\end{remark}

\subsection{马太效应与中美两极}

SCX审计基础设施不会均匀地分布在全球数据经济中。如同金融中心（纽约、伦敦、香港）和互联网交换点（阿姆斯特丹、法兰克福、圣何塞）的空间集中，数据审计基础设施必然在少数节点上汇聚。但SCX的集中逻辑不同于传统的网络效应或集聚经济——它是一种审计驱动的累积因果机制，其动力学由Theorem 3（SCX不确定性原理）的一个被忽视的推论所支配。

\begin{definition}[审计驱动的数据质量正反馈循环]
设经济体 $i$ 在时间 $t$ 拥有数据池 $\mathcal{D}_i^t$ 和模型质量 $M_i^t$。SCX审计下的动态演化由以下联立方程描述：
\begin{align}
Q_i^t &= \theta(|\mathcal{E}_i^t|) \cdot \text{Audit}(\mathcal{D}_i^t) \label{eq:quality}\\
M_i^{t+1} &= f(M_i^t, Q_i^t, C_i) \label{eq:model}\\
U_i^{t+1} &= g(M_i^{t+1}) \label{eq:users}\\
|\mathcal{E}_i^{t+1}| &= |\mathcal{E}_i^t| + h(U_i^{t+1}) \label{eq:audit-growth}
\end{align}
其中 $Q_i^t$ 是经Yajie审计认证的数据质量，$\theta(\cdot)$ 是CEC驱动的审计精度函数（公式(3)），$C_i$ 是算力禀赋，$U_i^{t+1}$ 是用户/采用者规模，$h(\cdot)$ 是将用户规模映射为审计数据增量的函数。
\end{definition}

方程(\ref{eq:quality})–(\ref{eq:audit-growth})描述了一个正反馈循环：更高的审计精度 → 更高的认证数据质量 → 更好的模型性能 → 更多的用户 → 更多的审计数据 → CEC扩展 → 审计精度进一步升高。这是一个典型的累积因果（cumulative causation）过程，与Myrdal (1957) 和 Arthur (1990) 所分析的区域发展极化机制在形式结构上同构。

\textbf{Theorem 3与不可区分性困境。} 这一反馈循环对后发者的影响并非仅仅是追赶困难——而是更根本的\textit{不可区分性}问题。考虑一个后发经济体 $j$，它在时间 $T$ 开始开发自己的SCX审计基础设施和AI模型。在任意时间 $t > T$，经济体 $j$ 观察到自己的模型性能 $M_j^t$ 低于先发者 $i$ 的性能 $M_i^t$。经济体 $j$ 面临一个诊断问题：$M_j^t < M_i^t$ 的原因是 (a) $j$ 的数据质量 $\mathcal{D}_j^t$ 低于 $i$，还是 (b) $j$ 的模型算法 $f_j$ 弱于 $i$ 的 $f_i$，还是 (c) 两者的某种混合？

Theorem 3（SCX不确定性原理）的一个直接推论是：在没有独立的、经过充分校准的审计协议的情况下，数据集的质量无法被其所有者可靠地自评。后发者的数据尚未通过Yajie级别的审计——因为Yajie的校准精度是CEC规模的函数，而后发者的CEC远小于先发者——因此后发者对其自身数据质量的估计本身就带有系统性的不确定性。这个不确定性与模型性能的因果归因不确定性相互缠绕，产生了一个不可分解的诊断困境：

\begin{equation}
\underbrace{M_j^t < M_i^t}_{\text{observable}} \quad \text{可归因于} \quad \begin{cases}
\mathcal{D}_j^t \text{ 质量较低} &\text{(数据假说)}\\
f_j \text{ 算法较弱} &\text{(算法假说)}\\
\text{两者的某种交互} &\text{(混合假说)}
\end{cases}
\end{equation}

而数据假说本身又依赖于一个审计基础设施的质量——CEC规模——来验证，这正是后发者所缺乏的。这是一个认知锁定的恶性循环：后发者需要高质量的审计来诊断自己的问题，但高质量的审计只能来自大规模的审计历史累积，而审计历史累积又需要先有一个确定性的、可信的数据质量基准来吸引初始用户。先发者在后发者进入市场之前，已经在无竞争对手的环境中完成了这一循环。

\begin{remark}[中美两极的必然性]
在SCX时代，两个经济体具备成为审计基础设施极点的全部必要条件：(1) 大规模、多样化的数据池（训练数据和审计校准数据）；(2) 独立的计算基础设施（GPU集群、数据中心、芯片制造能力）；(3) 足够大的内部市场以支持审计生态系统的初始积累（无需依赖外部用户即可达到CEC临界规模 $|\mathcal{E}|^*$）；(4) 数据治理的法律和制度框架，尽管内容不同（中国的数据安全法与GDPR代表了两种不同的治理哲学），但均具备将审计输出转化为监管合规证据的操作能力。

中国和美国各自满足这四项条件。欧盟在条件(2)（计算主权）和条件(3)（内部市场碎片化于27个成员国的监管体系之间）上存在结构性不足。印度在条件(1)（数据多样性的深度）和条件(2)（芯片制造）上处于追赶阶段。其他经济体在至少两个条件上存在不可忽视的缺口。

这意味着SCX时代的全球数据审计基础设施将以中美为两个天然极点。这不是一个规范性主张（normative claim），而是一个结构性推断（structural inference）：博弈的参数环境使得这两个经济体在审计驱动的累积因果循环中占据不可被后发竞争所侵蚀的初始位置。其他经济体——包括欧盟、日本、韩国、印度——必须作为审计接受者（audit receivers）参与这一体系：它们接入中美之一（或两者）的审计基础设施，将其数据质量认证委托给极点经济体维护的CEC校准锚。这不是殖民主义或依附关系——参与是自愿的，委托是理性的，审计输出是数学上可验证的——但其经济后果与依附理论（dependency theory）所描述的"中心–外围"结构（Prebisch, 1950; Frank, 1967）具有形式上的对应关系：中心生产审计能力，外围消费审计认证。
\end{remark}

\subsection{制造业复苏：审计的物理不可远程性}

前述分析中隐含着一个对全球经济地理具有深远影响的假设：数据可以从任何地方被审计。这一假设对于纯数字数据是正确的——一个存储在云端的训练数据集可以从任何具有网络连接的位置被下载、索引和审计。但这一假设在数据\textit{来源于物理生产过程}时不再成立。

\begin{definition}[物理嵌入数据的审计不可远程性]
设 $\mathcal{D}_{\text{phys}}$ 为一组物理嵌入数据（physically embedded data）——即，其生成过程依赖于物理传感器、测量设备、实验室仪器或生产线质量检测站的数据。$\mathcal{D}_{\text{phys}}$ 的审计质量 $\theta_{\text{audit}}$ 是以下因素的函数：
\begin{equation}
\theta_{\text{audit}} = \theta\left( \mathcal{A}, \mathcal{S}, \mathcal{H}, \mathcal{D}_{\text{phys}} \right)
\end{equation}
其中 $\mathcal{A}$ 是审计模块集合，$\mathcal{S}$ 是传感器网络的物理配置（空间分辨率、采样频率、校准状态），$\mathcal{H}$ 是现场领域专家的人为判断（异常解释、设备故障识别、环境因素校正）。$\mathcal{S}$ 和 $\mathcal{H}$ 不可通过远程通信完全复制——传感器需要物理安装和定期维护，领域专家需要对生产环境进行体感观察以形成可靠的异常判断。因此，物理嵌入数据的SCX审计是一个\textit{不可远程化}的过程：$\partial \theta_{\text{audit}} / \partial (\text{distance}) < 0$。
\end{definition}

这一性质的推论是直接而深远的：谁拥有生产线，并且在该生产线上部署了SCX审计基础设施（传感器+测量站+领域专家+Spring实例），谁就拥有了一种数学上可证明的、不可被竞争对手通过远程手段复制的质量优势。这个优势不是来自品牌声誉（可以被营销投资侵蚀），不是来自规模经济（可以被更大规模的竞争对手超越），也不是来自知识产权（可以被专利时效或反向工程削弱）。它来自一个数学事实：经SCX多专家共识验证的生产数据质量，为最终产品的品质提供了可独立核实、可在合同和监管中引用、可被下游客户作为采购决策依据的保证。

\begin{remark}[先发优势的转型：从廉价劳动力到数据质量]
制造业竞争力的历史叙事将先发优势定位于生产要素成本——特别是劳动力成本。英国工业革命的成功部分源于煤炭的廉价获取（Wrigley, 1988）；东亚奇迹（日本、韩国、中国）的早期阶段依赖于受过良好教育的、相对于其生产力水平而言廉价的劳动力（Amsden, 1989; Wade, 1990）。在这一传统叙事中，制造业随着劳动力成本的上升而自然迁移——从英国到美国，从美国到日本，从日本到亚洲四小龙，从亚洲四小龙到中国，以及正在发生的从中国到东南亚和南亚的迁移。

SCX审计基础设施切断了这一迁移逻辑。如果一个生产基地已经部署了SCX审计——物理传感器、校准测量站、经过审计训练的领域专家、与极点CEC锚定的Spring实例——那么它所生产的不仅是最终产品，还有一个不可分割的副产品：经数学验证的品质证明。将生产基地迁移到一个劳动力成本更低但缺乏SCX审计基础设施的地区，意味着失去这个品质证明，并因此失去其在全球市场中的质量溢价。节省的劳动力成本必须与以下损失进行比较：(a) 无法证明产品质量（导致价格折扣），(b) 重新积累审计历史的时间成本（CEC冷启动问题），(c) 在新地点重新部署物理传感器网络和培训领域专家的固定投资，以及 (d) 在新审计基础设施达到与旧设施相当的精度之前（可能需要数年时间）承受的竞争劣势。

在SCX审计覆盖面达到临界水平后，先发优势的参数从劳动力成本差异 $\Delta w$ 转移到审计质量差异 $\Delta \theta_{\text{audit}}$。一个在高劳动力成本地区运营但具有完备SCX审计的生产基地，在均衡中可能比一个在低劳动力成本地区运营但无审计的生产基地具有更高的利润率——因为市场为可验证的质量支付的价格溢价超过了劳动力成本的节省。
\end{remark}

\textbf{审计驱动的制造业回流。} 上述分析意味着一种新型的制造业回流（reshoring）机制——它不是由关税、贸易战或供应链安全的考虑驱动的（尽管这些因素也在起作用），而是由审计基础设施的物理不可远程性驱动的。一个在本国部署了SCX审计的制造企业，如果将生产外包到一个没有等效审计基础设施的海外地点，相当于放弃了其产品质量的数学可证明性。而一旦市场中的竞争对手（无论是国内的还是其他高审计密度地区的）向客户提供经Yajie认证的、可独立验证的品质数据，未审计的制造商就无法在质量维度上竞争。它们被迫在价格维度上竞争——而价格竞争在劳动力成本差异缩小的全球趋势中是边际递减的优势来源。

这一机制的净效应是：SCX审计基础设施的空间分布将成为制造业地理格局的一个决定性变量——其重要性不亚于劳动力成本、运输基础设施、能源价格或制度质量。拥有高密度SCX审计基础设施的国家（由于上述的马太效应，主要是中国和美国）将看到其制造业部门获得一种新型的比较优势：不是生产得更便宜，而是对生产质量的\textit{证明}更可信。不拥有这种审计基础设施的国家将面临一个选择：要么建设自己的审计基础设施（面临CEC冷启动和递归的不可区分性问题），要么将审计认证外包给极点经济体（接受审计接受者的角色），要么在无审计认证的状态下竞争（承受质量折扣）。

\section*{References}

\bibitem{adalovelace2020} Ada Lovelace Institute. \textit{Examining the Black Box: Tools for Assessing Algorithmic Systems}. London (2020).

\bibitem{armstrong2006} Armstrong, M. Competition in two-sided markets. \textit{The RAND Journal of Economics} \textbf{37}(3), 668--691 (2006).

\bibitem{arthur1989} Arthur, W. B. Competing technologies, increasing returns, and lock-in by historical events. \textit{The Economic Journal} \textbf{99}(394), 116--131 (1989).

\bibitem{bateman2022} Bateman, J. \textit{U.S.-China Technological ``Decoupling'': A Strategy and Policy Framework}. Carnegie Endowment (2022).

\bibitem{batini2009} Batini, C., Cappiello, C., Francalanci, C. \& Maurino, A. Methodologies for data quality assessment and improvement. \textit{ACM Computing Surveys} \textbf{41}(3), 1--52 (2009).

\bibitem{david1985} David, P. A. Clio and the economics of QWERTY. \textit{The American Economic Review} \textbf{75}(2), 332--337 (1985).

\bibitem{diakopoulos2015} Diakopoulos, N. Algorithmic accountability. \textit{Digital Journalism} \textbf{3}(3), 398--415 (2015).

\bibitem{drezner2015} Drezner, D. W. Targeted sanctions in a world of global finance. \textit{International Interactions} \textbf{41}(4), 755--764 (2015).

\bibitem{ec2021} European Commission. \textit{Proposal for a Regulation Laying Down Harmonised Rules on Artificial Intelligence}. COM(2021) 206 final (2021).

\bibitem{farrell2019} Farrell, H. \& Newman, A. L. Weaponized interdependence. \textit{International Security} \textbf{44}(1), 42--79 (2019).

\bibitem{farrell2007} Farrell, J. \& Klemperer, P. Coordination and lock-in. In \textit{Handbook of Industrial Organization}, Vol.~3, pp.~1967--2072. Elsevier (2007).

\bibitem{farrell1985} Farrell, J. \& Saloner, G. Standardization, compatibility, and innovation. \textit{The RAND Journal of Economics} \textbf{16}(1), 70--83 (1985).

\bibitem{floridi2018} Floridi, L., Cowls, J., Beltrametti, M., et al. AI4People---An ethical framework for a good AI society. \textit{Minds and Machines} \textbf{28}(4), 689--707 (2018).

\bibitem{gebru2021} Gebru, T., Morgenstern, J., Vecchione, B., Vaughan, J. W., Wallach, H., Daum\'e III, H. \& Crawford, K. Datasheets for datasets. \textit{Communications of the ACM} \textbf{64}(12), 86--92 (2021).

\bibitem{gregory2021} Gregory, R. W., Henfridsson, O., Kaganer, E. \& Kyriakou, H. The role of artificial intelligence and data network effects. \textit{Academy of Management Review} \textbf{46}(3), 534--551 (2021).

\bibitem{jobin2019} Jobin, A., Ienca, M. \& Vayena, E. The global landscape of AI ethics guidelines. \textit{Nature Machine Intelligence} \textbf{1}(9), 389--399 (2019).

\bibitem{joyner2011} Joyner, D. H. \textit{Interpreting the Nuclear Non-Proliferation Treaty}. Oxford University Press (2011).

\bibitem{kaplan2020} Kaplan, J., McCandlish, S., Henighan, T., Brown, T. B., Chess, B., Child, R., Gray, S., Radford, A., Wu, J. \& Amodei, D. Scaling laws for neural language models. \textit{arXiv preprint arXiv:2001.08361} (2020).

\bibitem{katz1985} Katz, M. L. \& Shapiro, C. Network externalities, competition, and compatibility. \textit{The American Economic Review} \textbf{75}(3), 424--440 (1985).

\bibitem{khan2017} Khan, L. M. Amazon's antitrust paradox. \textit{Yale Law Journal} \textbf{126}(3), 710--805 (2017).

\bibitem{mearsheimer1993} Mearsheimer, J. J. The case for a Ukrainian nuclear deterrent. \textit{Foreign Affairs} \textbf{72}(3), 50--66 (1993).

\bibitem{mokander2021} M\"{o}kander, J., Morley, J., Taddeo, M. \& Floridi, L. Ethics-based auditing of automated decision-making systems. \textit{Science and Engineering Ethics} \textbf{27}(4), 1--30 (2021).

\bibitem{nist2023} NIST. \textit{Artificial Intelligence Risk Management Framework 1.0}. U.S. Department of Commerce (2023).

\bibitem{northcutt2021} Northcutt, C. G., Athalye, A. \& Mueller, J. Pervasive label errors in test sets destabilize machine learning benchmarks. \textit{Proc. NeurIPS} (2021).

\bibitem{parasuraman1997} Parasuraman, R. \& Riley, V. Humans and automation: Use, misuse, disuse, abuse. \textit{Human Factors} \textbf{39}(2), 230--253 (1997).

\bibitem{parker2005} Parker, G. G. \& Van Alstyne, M. W. Two-sided network effects. \textit{Management Science} \textbf{51}(10), 1494--1504 (2005).

\bibitem{pushkarna2022} Pushkarna, M., Zaldivar, A. \& Kjartansson, O. Data cards: Purposeful and transparent dataset documentation. \textit{Proc. FAccT} 1776--1826 (2022).

\bibitem{raji2020} Raji, I. D., Smart, A., White, R. N., Mitchell, M., Gebru, T., Hutchinson, B., Smith-Loud, J., Theron, D. \& Barnes, P. Closing the AI accountability gap. \textit{Proc. FAccT} 33--44 (2020).

\bibitem{rochet2003} Rochet, J.-C. \& Tirole, J. Platform competition in two-sided markets. \textit{Journal of the European Economic Association} \textbf{1}(4), 990--1029 (2003).

\bibitem{sagan1996} Sagan, S. D. Why do states build nuclear weapons? \textit{International Security} \textbf{21}(3), 54--86 (1996).

\bibitem{sambasivan2021} Sambasivan, N., Kapania, S., Highfill, H., Akrong, D., Paritosh, P. \& Aroyo, L. M. ``Everyone wants to do the model work, not the data work'': Data cascades in high-stakes AI. \textit{Proc. CHI} 1--15 (2021).

\bibitem{sandvig2014} Sandvig, C., Hamilton, K., Karahalios, K. \& Langbort, C. Auditing algorithms: Research methods for detecting discrimination on internet platforms. \textit{Data and Discrimination: Converting Critical Concerns into Productive Inquiry} 1--23 (2014).

\bibitem{segal2020} Segal, A. The coming tech Cold War with China. \textit{Foreign Affairs} \textbf{99}(5), 94--105 (2020).

\bibitem{waltz1981} Waltz, K. N. The spread of nuclear weapons: More may be better. \textit{The Adelphi Papers} \textbf{21}(171), 1--32 (1981).

\bibitem{wang1996} Wang, R. Y. \& Strong, D. M. Beyond accuracy: What data quality means to data consumers. \textit{Journal of Management Information Systems} \textbf{12}(4), 5--33 (1996).

\bibitem{wilson2021} Wilson, C., Ghosh, A., Jiang, S., Mislove, A., Baker, L., Szurley, J., Dud\'ik, M., Danks, D. \& Ramachandran, P. Building and auditing fair algorithms. \textit{Proc. FAccT} 666--677 (2021).

\bibitem{zeff2012} Zeff, S. A. The evolution of the IASC into the IASB. \textit{The Accounting Review} \textbf{87}(3), 807--837 (2012).

\bibitem{zhang2014} Zhang, Y. \& Andrew, J. Financialisation and the conceptual framework. \textit{Critical Perspectives on Accounting} \textbf{25}(1), 17--26 (2014).

\end{thebibliography}

\end{document}
